
\documentclass[10pt]{article}
\usepackage[T1]{fontenc}
\usepackage[utf8]{inputenc}
\usepackage{lmodern}
\usepackage{amsmath,amssymb,amsthm,mathtools}
\usepackage{geometry}
\usepackage{enumitem}
\geometry{margin=0.86in}
\setlength{\parindent}{1.35em}
\setlength{\parskip}{0.18em}
\emergencystretch=3em
\setlist{leftmargin=2.2em,itemsep=0.12em,topsep=0.25em}
\newcommand{\A}{\mathbb A}
\newcommand{\Af}{\mathbb A^f}
\newcommand{\C}{\mathbb C}
\newcommand{\Gm}{\mathbb G_m}
\newcommand{\Q}{\mathbb Q}
\newcommand{\R}{\mathbb R}
\newcommand{\Z}{\mathbb Z}
\newcommand{\bS}{\mathbb S}
\newcommand{\GL}{\operatorname{GL}}
\newcommand{\Gp}{\operatorname{Gp}}
\newcommand{\Hom}{\operatorname{Hom}}
\newcommand{\Isom}{\operatorname{Isom}}
\newcommand{\Lie}{\operatorname{Lie}}
\newcommand{\Gal}{\operatorname{Gal}}
\newcommand{\Ker}{\operatorname{Ker}}
\newcommand{\ad}{\operatorname{ad}}
\newcommand{\dfn}{\operatorname{dfn}}
\newcommand{\id}{\operatorname{Id}}
\newcommand{\limproj}{\varprojlim}
\newcommand{\limind}{\varinjlim}
\newcommand{\MC}{M_{\C}}
\newcommand{\KMC}{K M_{\C}}
\newcommand{\whZ}{\widehat{\Z}}
\newcommand{\pii}{\pi_0}
\newcommand{\HasseH}{H^1}

\newcommand{\Spec}{\operatorname{Spec}}
\newcommand{\Res}{\operatorname{Res}}
\newcommand{\Norm}{\operatorname{N}}
\newcommand{\ab}{\mathrm{ab}}

\newcommand{\End}{\operatorname{End}}
\newcommand{\Tr}{\operatorname{Tr}}
\newcommand{\NS}{\operatorname{NS}}
\newcommand{\Aut}{\operatorname{Aut}}
\newcommand{\Id}{\operatorname{Id}}


\newcommand{\Sp}{\operatorname{Sp}}
\newcommand{\codim}{\operatorname{codim}}

\begin{document}


\begin{center}
Séminaire BOURBAKI\hfill Février 1971\\
23e année, 1970/71, n\textsuperscript{o} 389\\[2.2em]
{\Large\bfseries TRAVAUX DE SHIMURA}\footnote{Texte rédigé en mars 1971}\\[1.2em]
par Pierre DELIGNE
\end{center}

Soit $X/\Gamma$ le quotient d'un domaine hermitien symétrique $X$ par un groupe discret arithmétique $\Gamma$ (supposé défini par des conditions de congruence, voir 1.7). D'après Baily et Borel, $X/\Gamma$ est une variété algébrique complexe. Dans les articles [10] à [14], Shimura montre entre autres, pour beaucoup de ces variétés, qu'elles peuvent être définies sur des corps de nombres qu'il explicite. C'est là le sujet du présent exposé.

Pour obtenir des énoncés propres, on est contraint à utiliser un langage adélique. Celui-ci conduit à considérer, pour une ``condition de congruence'' donnée, un schéma sur $\C$, somme disjointe d'un nombre fini de variétés du type $X/\Gamma$ (voir §§ 1 et 2). Dans de nombreux cas, ce schéma pourra être défini sur un corps de nombres $E$, indépendant de la condition de congruence considérée. Ses composantes connexes géométriques, elles, seront définies sur des corps de classe de $E$. Ce phénomène, bien qu'essentiel, va être négligé dans la suite de cette introduction.

Dans un petit nombre de cas, $X/\Gamma$ peut s'interpréter comme l'ensemble des classes d'isomorphie des variétés abéliennes complexes, munies de quelques structures algébriques additionnelles (polarisations, endomorphismes, structures sur les points d'ordre $n$). La solution, sur $\Q$, d'un problème de modules correspondant fournit alors un schéma $M$ sur un corps de nombres explicite $F$, dont $X/\Gamma$ est l'ensemble des points complexes. On appellera $M$ un ``modèle'' de $X/\Gamma$. Le cas fondamental est celui des variétés abéliennes polarisées de la série principales, munies d'une structure de niveau $n$ (1.6, 1.11, 4.16, 4.17 et 4.21).

Dans d'autres cas, $X/\Gamma$ peut s'interpréter comme l'ensemble des classes d'isomorphie des variétés abéliennes complexes $A$ munies de quelques structures comme plus haut, et de quelques classes de cohomologie entières de type $(p,p)$
\[
a_i\in H^{2p}(A\times\cdots\times A,\C) \qquad (\text{cf. Mumford [7]}).
\]
Dans ce cas, l'idée sera de construire un modèle $M$ de $X/\Gamma$ comme sous-schéma d'un modèle $M'$ déjà construit d'une variété $X'/\Gamma'$. En gros, on commence par construire dans $M'$ les points de $M$ correspondant à des variétés abéliennes de type C.M. (= maximum de multiplication complexe). On définit alors $M$ comme l'adhérence de l'ensemble de ces points. Cette construction repose sur la connaissance détaillée des variétés abéliennes de type C.M. fournie par [16].

Dans les cas considérés plus haut, on dispose de beaucoup d'information sur les corps de définition des points de $M$ correspondant à des variétés abéliennes de type C.M. Ceci permet de donner des solutions partielles au 12e problème de Hilbert (cf. l'introduction de [15] et 3.15, 3.16).

En regardant ce que ces cas ont en commun, on aboutit à la notion de ``modèle canonique'' (§ 3). On montre qu'il y a au plus un ``modèle canonique'' des $X/\Gamma$ (5.5). Des techniques de descente permettent alors de construire quelques modèles canoniques qui semblent loin d'être des solutions de problèmes de modules ([12] [13] [14] et [5]). Pour la description de ces ``modèles étranges'', on renvoie aux §§ 5 et 6.

On désignera par $\Z/n(1)$ tant le schéma en groupe $\mu_n$ des racines $n$-ièmes de l'unité que le groupe des racines $n$-ièmes de l'unité dans un corps algébriquement clos fixé $\bar{k}$. Pour $\bar{k}=\C$, l'exponentielle permet d'identifier $\Z/n(1)$ et $\Z/n$. Les $\Z/n(1)$ forment un système projectif : les applications de transition sont les
\[
\varphi_{n,m}:x\longmapsto x^{n/m}.
\]
On pose
\[
\whZ=\limproj_n \Z/n\Z \simeq \prod_p \Z_p,\qquad
\whZ(1)=\limproj_n \Z/n\Z(1),
\]
\[
\Af=\Q\otimes\whZ\simeq \prod_p{}'\Q_p\quad(\text{produit restreint}),\qquad
\Af(1)=\Q\otimes\whZ(1)=\Af\otimes_{\whZ}\whZ(1),
\]
et on désigne par $\A$ l'anneau des adèles de $\Q$.

On utilisera les notations suivantes :

\noindent $\pii(X)$ (pour $X$ un espace topologique) : l'ensemble des composantes connexes de $X$. Dans la suite, $\pii(X)$, muni de la topologie quotient de celle de $X$, sera discret ou compact totalement discontinu.

\noindent $G^0$ : composante connexe de l'origine d'un groupe topologique $G$ (ou d'un espace topologique pointé).

\noindent $G(K)$, $G_K$, $G\otimes_FK$ : pour $G$ un schéma sur $F$ (par exemple un groupe algébrique sur $F$) et $K$ une $F$-algèbre, on désigne par $G(K)$ l'ensemble des points de $G$ à valeurs dans $K$ et par $G_K$ ou par $G\otimes_FK$ le schéma sur $K$ déduit de $G$ par extension des scalaires.

\noindent $E^*$ : pour $E$ une algèbre de rang fini sur un corps $F$, $E^*$ désigne le groupe algébrique sur $F$ des éléments inversibles de $E$. Pour $E$ un corps de nombres et $F=\Q$, $E^*(\A)/E^*(\Q)$ est le groupe des classes d'idèles de $E$.

On dira qu'un groupe algébrique $G$ sur $\Q$ vérifie le principe de Hasse si
\[
H^1(\Q,G)=\dfn H^1(\Gal(\overline\Q/\Q),G(\overline\Q))
\]
s'injecte dans le produit étendu à toutes les places des $H^1(\Q_v,G_{\Q_v})$. On utilisera les théorèmes suivants.

\paragraph{(0.1) Principe de Hasse.} Un groupe semi-simple simplement connexe, sans facteur de type $E_8$, vérifie le principe de Hasse [3].

\paragraph{(0.2)} Un groupe semi-simple simplement connexe $G$ sur un corps local non archimédien $K$ vérifie $H^1(K,G)=0$ (une démonstration indépendante de la classification est annoncée dans Bruhat--Tits [1]).

\paragraph{(0.3) Théorème d'approximation fort.} Soit $G$ un groupe semi-simple simplement connexe sur un corps global $K$. Si $G$ est $K$-simple, et si $v$ est une place de $K$ telle que $G(K_v)$ soit non compact, alors $G(K_v)G(K)$ est dense dans $G(\A)$ (pour une démonstration indépendante de la classification, voir Platonov [9]).

\paragraph{(0.4) Théorème d'approximation réel.} Soit $G$ un groupe algébrique connexe (linéaire) sur $\Q$. Alors $G(\Q)$ est dense dans $G(\R)$ (pour le prouver, on se ramène aisément au cas des tores, qui est un corollaire de ([4] 5.1)).

\section*{§ 1. Le langage adélique.}

\paragraph{1.1.} Soit $G$ un groupe réductif sur $\Q$. On désigne par $G'$ le groupe dérivé de la composante neutre de $G$, par $C$ le centre de cette composante neutre, et on pose $T=G/G'$. On dispose de suites exactes de groupes algébriques
\begin{align}
0&\longrightarrow C\longrightarrow G\longrightarrow G/C\longrightarrow0,\tag{1.1.1}\\
0&\longrightarrow G'\longrightarrow G\xrightarrow{\nu}T\longrightarrow0,\tag{1.1.2}
\end{align}
et les applications canoniques de $C$ dans $T$ et de $G'$ dans $G/C$ sont des isogénies de noyau $C\cap G'$.

\paragraph{1.2.} Soit $\bS$ le groupe algébrique réel des éléments inversibles de la $\R$-algèbre $\C$. C'est encore le groupe algébrique réel déduit par restriction des scalaires à la Weil, de $\C$ à $\R$, du groupe multiplicatif $\Gm$. On a $\bS(\R)=\C^*$, et $\bS_\C\simeq\mathbb G_{m,\C}\times\mathbb G_{m,\C}$. Le groupe $\Hom(\bS_\C,\mathbb G_{m,\C})$ des caractères de $\bS_\C$ a pour base les caractères $z$ et $\bar z$ tels que le composé
\[
\C^*=\bS(\R)\longrightarrow \bS(\C)\xrightarrow{z,\bar z}\C^*
\]
soit respectivement l'identité et la conjugaison complexe.

Soit $V$ un espace vectoriel réel. Pour toute représentation $h:\bS\to\GL(V)$, on désigne par $V^{pq}$ le sous-espace de $V_\C$ sur lequel $\bS$ agit par le caractère $z^p\bar z^q$. Les $V^{pq}$ forment une bigraduation de $V_\C$, et la construction $h\mapsto V^{pq}$ identifie les représentations de $\bS$ dans $\GL(V)$ aux bigraduations de Hodge de $V_\C$, i.e. aux bigraduations de $V_\C$ telles que $V^{pq}=\overline{V^{qp}}$. On désignera par $F(h)$ la filtration de Hodge de $V_\C$ :
\[
F(h)^p=\bigoplus_{p'\ge p}V^{p'q'}.
\]

\paragraph{1.3.} On désigne par $r:\mathbb G_{m,\C}\to\bS_\C$ le $\C$-homomorphisme tel que
\[
(z^p\bar z^q)\circ r=(x\longmapsto x^p),
\]
et par $w:\mathbb G_{m,\R}\to\bS$ le $\R$-homomorphisme tel que
\[
(z^p\bar z^q)\circ w=(x\longmapsto x^{p+q}).
\]
Pour $h:\bS\to\GL(V)$ une représentation et $v^{pq}\in V^{pq}$, on a
\[
hr(x)\,v^{pq}=x^p v^{pq}\quad\text{et}\quad hw(x)\,v^{pq}=x^{p+q}v^{pq}.
\]
Le composé $w:\R^*=\mathbb G_m(\R)\xrightarrow{w}\bS(\R)=\C^*$ est l'inclusion évidente.

\paragraph{1.4.} Soit $h:\bS\to G_\R$ un homomorphisme. Pour toute représentation $\rho:G\to\GL(V)$ de $G$ sur un $\Q$-espace vectoriel $V$, $\rho h$ définit une bigraduation de $V_\C$ (1.2). Si $V$ est fidèle, et si $G$ est le sous-groupe de $\GL(V)$ laissant fixes quelques tenseurs $s_i$, la construction $h\mapsto(V^{pq})$ identifie les homomorphismes $h$ aux bigraduations de Hodge de $V_\C$ telles que les $s_i$ soient de type $(0,0)$.

\paragraph{1.5.} Soit $h_o:\bS\to G_\R$ un homomorphisme vérifiant les conditions suivantes (cf. [7] et [2]; (1.5.2) est motivé par le théorème de transversalité de Griffiths).

\begin{enumerate}[label={(1.5.\arabic*)},leftmargin=3.0em]
\item L'image de $h_ow:\Gm\to G_\R$ est centrale. On appellera $h_ow$ le poids de $h_o$;
\item la bigraduation de Hodge du complexifié $\Lie(G)_\C$ de l'espace de la représentation adjointe (1.2) est de type $(-1,1)+(0,0)+(1,-1)$;
\item l'automorphisme $\ad h_o(i)$ de $G$, une involution d'après (1.5.1), induit une involution de Cartan de $G'$.
\end{enumerate}

Soit $K_\infty$ le centralisateur de $h_o$ dans $G(\R)$; il contient le centre de $G(\R)$, et $K_\infty\cap G'(\R)^0$ est un sous-groupe compact maximal de $G'(\R)^0$, identique au centralisateur dans $G'(\R)^0$ de $h_o(i)$. L'espace $K_\infty\backslash G(\R)$ s'identifie à l'ensemble $X$ des conjugués de $h_o$ par un élément de $G(\R)$. On vérifie qu'il possède une et une seule structure complexe telle que, pour toute représentation $\rho:G\to\GL(V)$, la filtration $F(\rho h)$ de $V_\C$ dépende holomorphiquement de $h\in X$. Cette structure est invariante à droite sous $G(\R)$, et les composantes connexes de $X$ sont des domaines hermitiens symétriques. On désignera par $X^0$ la composante connexe de $h_o$.

\paragraph{1.6. Exemple I.} Soit $\Gp(V)$ le groupe des similitudes symplectiques d'un vectoriel réel $V$ muni d'une forme alternée non dégénérée $\psi$. Il existe alors une et une seule classe de conjugaison $X$ d'homomorphismes $h:\bS\to\Gp(V)$ vérifiant les conditions (1.5.1) et de poids $-1$, i.e. de poids
\[
hw:\Gm\to\Gp(V):x\longmapsto\text{homothétie de rapport }x^{-1}.
\]
L'espace $X$ s'identifie à la somme de deux demi-espaces de Siegel.

La construction 1.5 identifie les $h\in X$ aux décompositions
\[
V_\C=V^{-1,0}\oplus V^{0,-1}
\]
de $V_\C=V\otimes_\R\C$, telles que $V^{pq}=\overline{V^{qp}}$, que $V^{0,-1}$ soit totalement isotrope pour $\psi$, et que, si $C$ désigne l'endomorphisme de $V$ qui induit la multiplication par $i^{p-q}$ sur $V^{pq}$, la forme symétrique $\psi(x,Cy)=\psi(x,h(i)y)$ sur $V$ soit définie ($>0$ ou $<0$).

\paragraph{1.7.} Soit $\Gamma$ un sous-groupe arithmétique de $G(\Q)$. On s'intéresse au cas où $\Gamma$ est défini par des conditions de congruence. Si $V_\Z$ est un réseau entier dans une représentation fidèle $V$ de $G$, ceci signifie qu'il existe $n$ tel que $\Gamma$ contienne avec un indice fini le groupe $\Gamma_n\subset G(\Q)$ des $\gamma\in G(\Q)$ tels que $\gamma\in G(\R)^0$, que $\gamma V_\Z=V_\Z$ et que $\gamma\equiv\id\pmod n$. D'après Baily et Borel, le quotient $X^0/\Gamma$ est muni d'une structure naturelle de schéma quasi-projectif sur $\C$. Soit $K_n$ le sous-groupe compact ouvert de $G(\Af)$ formé des $k=(k_p)$ tels que
\[
k_p\cdot V_\Z\otimes\Z_p=V_\Z\otimes\Z_p
\quad\text{et}\quad
k_p\equiv1\pmod n\quad\text{pour }p\mid n.
\]
Le groupe $\Gamma_n$ est alors l'intersection, dans $G(\A)$, de $G(\Q)$ et de $G(\R)^0\times K_n$.

\paragraph{1.8.} Soit $K$ un sous-groupe compact ouvert de $G(\Af)$. L'espace
\[
K_\infty\times K\backslash G(\A)/G(\Q)
=K\backslash X\times G(\Af)/G(\Q)
=X\times(K\backslash G(\Af))/G(\Q)
\]
est alors somme d'une famille finie d'espaces du type considéré en 1.7. C'est donc un schéma quasi-projectif sur $\C$. On le notera ${}_K\MC(G,h_o)$ (voire simplement ${}_K\MC(G)$ ou ${}_K\MC$).

Pour $K$ de plus en plus petit, ces schémas forment un système projectif de schémas à morphismes de transition finis et surjectifs
\[
{}_K\MC(G,h_o)\qquad(K\to0). \tag{1.8.1}
\]
On désignera par $\MC(G,h_o)$ (voire simplement $\MC(G)$, ou $\MC$) le schéma quasi-compact et séparé (non de type fini pour $G\ne\{1\}$) limite projective de (1.8.1)
\[
\MC(G,h_o)=\limproj_K K_\infty\times K\backslash G(\A)/G(\Q). \tag{1.8.2}
\]
On doit le considérer comme un avatar du système projectif (1.8.1). Le groupe $G(\Af)$ agit sur $\MC$, ou, si l'on préfère, sur le pro-objet défini par le système projectif (1.8.1), mais bien sûr pas sur les ${}_K\MC$ individuels (cf. 3.2). On a
\[
K\backslash\MC={}_K\MC,
\]
de sorte que
\[
\MC=\limproj_K K\backslash\MC\qquad(K\text{ compact ouvert dans }G(\Af)). \tag{1.8.3}
\]
Ceci s'exprime en disant que $G(\Af)$ agit continûment sur $\MC$.

La construction suivante sera utile pour interpréter les ${}_K\MC$ comme des schémas de modules.

\paragraph{1.9.} Soient $H$ un groupe algébrique sur $\Q$ et $V$ une représentation fidèle de $H$. On définit une $H$-structure $\bar\iota$ sur un vectoriel $W$ comme étant un isomorphisme de $V$ avec $W$, donné modulo $H(\Q)$ : $\bar\iota\in\Isom(V,W)/H(\Q)$. Si $H(\Q)$ est le groupe des automorphismes d'une structure $s$ d'espèce $\Sigma$ sur $V$, alors une $H$-structure sur $W$ s'identifie à une structure $t$ de même espèce sur $W$, isomorphe à $s$ (avec $\bar\iota=\Isom((V,s),(W,t))$). Soit $W$ muni d'une $H$-structure $\bar\iota$. Le groupe algébrique $\iota H\iota^{-1}\subset\GL(W)$ ($\iota\in\bar\iota$) ne dépend que de $\bar\iota$. On le note $H^{\bar\iota}$. Soit $A$ une $\Q$-algèbre. L'ensemble $\Isom_{\bar\iota}(W\otimes A,V\otimes A)$ des isomorphismes permis de $W\otimes A$ avec $V\otimes A$ est l'ensemble des isomorphismes $k$ tels que, pour $\iota\in\bar\iota$, on ait $k\iota\in H(A)$.

\paragraph{1.10.} Soient $(G,h_o)$ comme en 1.5, $V$ une représentation linéaire fidèle de $G$, $K$ un sous-groupe compact ouvert de $G(\Af)$, et considérons les objets $H$ du type suivant : $H$ consiste en
\begin{enumerate}[label=(\alph*),leftmargin=2.6em]
\item un espace vectoriel $H_\Q$ sur $\Q$, muni d'une $G$-structure $\bar\iota$ (1.9);
\item un homomorphisme $h:\bS\to G^{\bar\iota}_\R$ tel que, pour $\iota\in\bar\iota$, $\iota^{-1}h\iota:\bS\to G_\R$ soit conjugué à $h_o$;
\item une classe $\bar k\in K\backslash\Isom_{\bar\iota}(H\otimes\Af,V\otimes\Af)$ (1.9).
\end{enumerate}

\paragraph{1.11. Exemple I (suite de 1.6).} Soient $V$ un $\Q$-vectoriel muni d'une forme alternée non dégénérée $\psi$, $G$ le groupe des similitudes symplectiques de $V$, $V_\Z$ un réseau entier sur lequel $\psi$ soit à valeurs entières de discriminant $1$ et $h_o$ comme en 1.6. Prenons
\[
K=\prod_\ell \Gp(V_\Z\otimes\Z_\ell).
\]
Une donnée (a) s'interprète comme un vectoriel $H$ de même dimension que $V$, muni d'une forme alternée $\psi$ donnée à un facteur rationnel près. Une donnée (b) s'interprète comme une bigraduation de Hodge de $H_\C$, de type 1.6. Une donnée (c) s'interprète comme un réseau entier $H_\Z$ de $H$, sur lequel un multiple rationnel de $\psi$ est à valeurs entières et de discriminant $1$ : $\bar k$ est l'ensemble des $k$ tels que, pour tout $\ell$, $V_\Z\otimes\Z_\ell=k_\ell(H_\Z\otimes\Z_\ell)$. On peut normaliser $\psi$ par les conditions d'être à valeurs entières de discriminant $1$ sur $H_\Z$, et de vérifier
\[
\psi(x,Cx)>0 \qquad (\text{cf. }1.6).
\]
Soient $n$ un entier, et $K_n$ le sous-groupe de $K$ formé des $k$ tels que $k\equiv1\pmod n$. Pour $K_n$, une donnée (c) s'interprète comme $H_\Z$ comme plus haut, plus une similitude symplectique
\[
H_\Z/nH_\Z\xrightarrow{\sim}V_\Z/nV_\Z.
\]

\paragraph{1.12.} Sous les hypothèses de 1.10, les points de ${}_K\MC(G,h_o)$ sont en correspondance biunivoque avec les classes d'isomorphie d'objets $H$ du type 1.10. Soit en effet
\[
H=(H_\Q,\bar\iota,h,\bar k).
\]
À $H$ muni de $\iota\in\bar\iota$, on associe le morphisme $\iota^{-1}h\iota:\bS\to G_\R$, élément de $X\simeq K_\infty\backslash G(\R)$, et $\bar k$ dans $K\backslash G(\Af)$, soit au total un élément de $K\backslash X\times G(\Af)$; à $H$ seul, on associe seulement un élément de
\[
{}_K\MC(G,h_o)=K\backslash X\times G(\Af)/G(\Q)=K_\infty\times K\backslash G(\A)/G(\Q),
\]
qui détermine univoquement la classe d'isomorphisme de $H$.

Dans certains cas (voir 4.11), se donner un objet $H$ comme plus haut revient à se donner une variété abélienne, munie de quelques structures supplémentaires. On obtient alors une interprétation de ${}_K\MC(G)$ (et de $\MC(G)$) comme étant le schéma de modules (grossier) de ce type de variétés abéliennes.

\paragraph{1.13.} Soient $G^i$ ($i=1,2$) deux groupes réductifs, $h^i:\bS\to G^i_\R$ des homomorphismes vérifiant les conditions de 1.5 et $X^i$ l'ensemble des conjugués de $h^i$ par un élément de $G^i(\R)$. Soient $G=G^1\times G^2$, $h:\bS\to G_\R$ le morphisme $(h^1,h^2)$ et $X=X^1\times X^2$ l'ensemble des conjugués de $h$. On définit de façon évidente un isomorphisme
\[
\MC(G^1,h^1)\times\MC(G^2,h^2)=\MC(G,h). \tag{1.13.1}
\]
Pour $K^i$ compact ouvert dans $G^i(\Af)$ et $K=K^1\times K^2$, on a de même
\[
{}_{K^1}\MC(G^1,h^1)\times{}_{K^2}\MC(G^2,h^2)={}_K\MC(G,h). \tag{1.13.2}
\]

\paragraph{1.14.} Soit $G^1,G^2,h^1,h^2$ comme en 1.13. On désignera par $u:(G^1,h^1)\to(G^2,h^2)$ un homomorphisme $u:G^1\to G^2$ tel que $uX^1\subset X^2$. Un tel homomorphisme définit
\[
u:\MC(G^1,h^1)\longrightarrow\MC(G^2,h^2). \tag{1.14.1}
\]
Pour $K^1$ compact ouvert dans $G^1(\Af)$, avec $u(K^1)\subset K^2$, il définit
\[
u(K^1,K^2):{}_{K^1}\MC(G^1,h^1)\longrightarrow{}_{K^2}\MC(G^2,h^2). \tag{1.14.2}
\]

\paragraph{Proposition 1.15.} Avec les notations de 1.14, supposons que $G^1$ soit un sous-groupe de $G^2$. Alors, pour tout sous-groupe compact ouvert $K^1$ de $G^1(\Af)$, il existe un sous-groupe compact ouvert $K^2\supset K^1$ de $G^2(\Af)$ tel que $u(K^1,K^2)$ identifie ${}_{K^1}\MC(G^1,h^1)$ à un sous-schéma fermé de ${}_{K^2}\MC(G^2,h^2)$.

Si les $K^i$ sont assez petits pour que l'application de $X^i\times(K^i\backslash G^i(\Af))$ dans ${}_{K^i}\MC(G^i,h^i)$ soit étale, alors l'application $u(K^1,K^2)$ est finie et non ramifiée. Il suffit de traiter ce cas, et de montrer que $u(K^1,K^2)$ est injectif pour $K^2\supset K^1$ assez petit. Le graphe de la relation d'équivalence
\[
u(K^1,K^2)(x)=u(K^1,K^2)(y)
\]
est un sous-schéma fermé de $({}_{K^1}\MC(G^1,h^1))^2$, qui décroît avec $K^2$, donc est stationnaire pour $K^2\supset K^1$ assez petit. Il suffit de montrer que
\[
u(K^1):{}_{K^1}\MC(G^1,h^1)\longrightarrow\limproj_{K^2\supset K^1}{}_{K^2}\MC(G^2,h^2)=\dfn{}_{K^1}\MC(G^2,h^2)
\]
est injectif. On a
\[
{}_{K^1}\MC(G^1,h^1)=K^1\backslash \MC(G^1,h^1),
\qquad
{}_{K^1}\MC(G^2,h^2)=K^1\backslash \MC(G^2,h^2),
\]
de sorte qu'il suffit de prouver la

\paragraph{Variante 1.15.1.} $u:\MC(G^1,h^1)\longrightarrow\MC(G^2,h^2)$ est injectif.



Cette assertion se déduit aisément des deux lemmes suivants.

\paragraph{Lemme 1.15.2.} Soit $U^i\subset C^i(\Q)$ le groupe des unités du centre de $G^i$, et soit
\[
G^i(\Q)^\wedge=\limproj_n\bigl(G^i(\Q)/(U^i)^n\bigr).
\]
Alors
\[
\MC(G^i,h^i)=K^i_\infty\backslash G^i(\A)/G^i(\Q)^\wedge .
\]
C'est un corollaire du théorème de Chevalley selon lequel les $(U^i)^n$ sont des sous-groupes de congruence de $U^i$.

\paragraph{Lemme 1.15.3.} L'application
\[
u:G^1(\Af)/G^1(\Q)^\wedge\longrightarrow G^2(\Af)/G^2(\Q)^\wedge
\]
est injective.

Par translation à gauche, il suffit de montrer que $u^{-1}(\{e\})=\{e\}$; on contemple le diagramme
\[
\begin{array}{ccccccccc}
0&\longrightarrow&G^1(\Q)^\wedge&\longrightarrow&G^2(\Q)^\wedge&\longrightarrow&(G^2/G^1)(\Q)^\wedge\\
&&\downarrow&&\downarrow&&\downarrow\!\\
0&\longrightarrow&G^1(\Af)&\longrightarrow&G^2(\Af)&\longrightarrow&(G^2/G^1)(\Af)\\
&&\downarrow&&\downarrow&&\\
&&G^1(\Af)/G^1(\Q)^\wedge&\longrightarrow&G^2(\Af)/G^2(\Q)^\wedge.&&
\end{array}
\]
où
\[
(G^2/G^1)(\Q)^\wedge=\limproj_n (G^2/G^1)(\Q)/(U^2/U^1)^n .
\]

\section*{§ 2. Composantes connexes.}

\paragraph{2.1.} Soit $G$ un groupe réductif connexe sur $\Q$. On reprend les notations 1.1 et on suppose que $G'$ n'admet pas de sous-groupe invariant non trivial $H$ (défini sur $\Q$) tel que $H(\R)$ soit compact. Cette hypothèse n'exclut pas que $G'_\R$ ait des facteurs compacts. On désigne par $\widetilde G'$ le revêtement simplement connexe du groupe dérivé $G'$.

\paragraph{Proposition 2.2.} (i) Le groupe $G(\Af)$ agit transitivement sur l'ensemble
\[
\pi_0(G(\A)/G(\Q))
\]
des composantes connexes de $G(\A)/G(\Q)$.

(ii) Le groupe $\widetilde G'(\A)$ agit trivialement sur $\pi_0(G(\A)/G(\Q))$.

(iii) Le quotient de $G(\A)$ par l'image de $\widetilde G'(\A)$ est un groupe commutatif.

L'assertion (i) résulte de (0.4) :
\[
G(\A)=G(\Af)G(\R)^0G(\Q).
\]
Puisque $\widetilde G'$ est simplement connexe, le groupe $\widetilde G'(\R)$ est connexe. D'après (0.3), $\widetilde G'(\R)\widetilde G'(\Q)$ est dense dans $\widetilde G'(\A)$, de sorte que $\widetilde G'(\A)/\widetilde G'(\Q)$ est connexe. Pour $x\in G(\A)$, on a $\widetilde G'(\A)x=x\widetilde G'(\A)$, de sorte que l'image de $\widetilde G'(\A)x$ dans $G(\A)/G(\Q)$ est une image de $\widetilde G'(\A)/\widetilde G'(\Q)$ : elle est connexe, et (ii) en résulte.

Soit $Z$ le centre de $\widetilde G'$. L'assertion (iii) résulte de ce que l'application «commutateur»
\[
(x,y)\longmapsto xyx^{-1}y^{-1}:G\times G\longrightarrow G
\]
admet une factorisation
\[
G\times G\longrightarrow G/C\times G/C\longleftarrow \widetilde G'/Z\times \widetilde G'/Z\longrightarrow \widetilde G'\longrightarrow G.
\]

\paragraph{Définition 2.3.} On désigne par $\pi(G)$ le quotient (commutatif) de $G(\A)$ (ou de $G(\Af)$) qui agit fidèlement sur $\pi_0(G(\A)/G(\Q))$.

En d'autres termes, $\pi(G)$ est $\pi_0(G(\A)/G(\Q))$, muni de sa structure de groupe commutatif «évidente». Rappelons le théorème suivant.

\paragraph{Théorème 2.4.} Sous les hypothèses de 2.1, si $G'$ est simplement connexe, alors l'homomorphisme canonique
\[
\pi(G)=\pi_0(G(\A)/G(\Q))\longrightarrow \pi_0(T(\A)/T(\Q)) \tag{2.4.1}
\]
est bijectif.

Pour prouver la surjectivité de (2.4.1), il suffit, d'après 2.2(i) appliqué à $T$, de montrer que $G(\Af)$ s'envoie sur $T(\Af)$. Le noyau $G'$ de $\nu:G\to T$ étant connexe, on sait que, pour presque tout $p$, $\nu(G(\Z_p))=T(\Z_p)$ (cette formule a un sens pour presque tout $p$, et est un corollaire du théorème de Lang appliqué à $G'$ réduit modulo $p$). Il suffit donc de prouver que, pour tout $p$,
\[
\nu(G(\Q_p))=T(\Q_p),
\]
ce qui résulte de (0.2) appliqué à $G'$.

Le théorème 2.4 équivaut à l'énoncé plus concret suivant.

\paragraph{Variante 2.5.} Pour tout sous-groupe compact ouvert $K$ de $G(\Af)$, l'application
\[
\nu:\pi_0(K\backslash G(\A)/G(\Q))\longrightarrow\pi_0(\nu(K)\backslash T(\A)/T(\Q)) \tag{2.5.1}
\]
est bijective.

En effet,
\[
\pi_0(G(\A)/G(\Q))=\limproj_K\pi_0(K\backslash G(\A)/G(\Q)),
\]
et de même pour $T$. Par ailleurs,
\[
\pi_0(K\backslash G(\A)/G(\Q))=G(\R)^0\backslash K\backslash G(\A)/G(\Q).
\]
Montrons que si $x,y\in G(\A)$ ont même image dans le membre de droite de (2.5.1), alors ils ont même image dans le membre de gauche. Une translation à gauche par $y^{-1}$, qui remplace $K$ par $y^{-1}Ky$, permet de supposer que $y=e$. Modifiant $x$ par un élément de $G(\R)^0\times K$, on peut supposer que $\nu(x)=\tau$ avec $\tau\in T(\Q)$. D'après le principe de Hasse (0.1) pour $G'$, l'équation $\nu(y)=\tau$ a une solution dans $G(\Q)$. Corrigeant $x$ par $y$, on peut supposer que $x\in G'(\A)$. D'après 2.2(ii), l'image de $x$ dans $G(\A)/G(\Q)$ est dans la composante neutre, et a fortiori $x$ et $y$ ont même image dans le membre de gauche de 2.5.1.

Cette démonstration semble utiliser, via (0.1), que $G'$ soit sans facteur $E_8$. En fait, la classe dans $H^1(\Q,G')$ qui obstrue l'équation $\nu(y)=\tau$ vient du centre de $G'$, et les facteurs $E_8$, étant adjoints, comptent pour du beurre.

\paragraph{Proposition 2.6.} On a
\[
K_\infty=C(\R)\bigl(K_\infty\cap G'(\R)^0\bigr).
\]
Le groupe $K_\infty\cap G'(\R)^0$ est connexe, et
\[
\pi_0(K_\infty)\xrightarrow{\sim}\operatorname{Im}\bigl(\pi_0(C(\R))\longrightarrow\pi_0(G(\R))\bigr).
\]

Soit $h':\bS\to G_\R\to(G/C)_\R$. Le centralisateur de $h'$ dans $(G/C)(\R)$ est connexe, car c'est un groupe compact dont le complexifié est connexe (comme centralisateur d'un tore). En particulier, l'image de $K_\infty$ dans $G/C(\R)$ tombe dans la composante neutre, image de $G'(\R)^0$, et
\[
K_\infty=C(\R)\bigl(K_\infty\cap G'(\R)^0\bigr).
\]
Le groupe $K_\infty\cap G'(\R)^0$ est connexe en tant que sous-groupe compact maximal d'un groupe connexe. 2.6 en résulte.

\paragraph{2.7.} L'image de $\pi_0(K_\infty)$ dans $\pi_0(G(\A)/G(\Q))$ ne dépend que de la classe de conjugaison de $h$, et
\[
\pi_0(\MC(G,h))=\limproj_K\pi_0(K_\infty\times K\backslash G(\A)/G(\Q))
\]
s'identifie à $\pi(G)/\pi_0(K_\infty)$.

En particulier, si $G'$ est simplement connexe, on a
\[
\pi_0(\MC(G,h))\xrightarrow{\sim}\pi_0(T(\A)/T(\Q))/\pi_0(K_\infty). \tag{2.7.1}
\]
Plus concrètement, pour tout sous-groupe compact ouvert $K$ de $G(\Af)$, $\pi_0({}_K\MC(G,h))$ est alors un espace principal homogène sous
\[
\nu(K_\infty\times K)\backslash T(\A)/T(\Q).
\]

\section*{§ 3. Modèles.}

Soient $G$ un groupe réductif sur $\Q$, $h:\bS\to G_\R$ vérifiant les conditions de 1.5, et $E$ un corps muni d'un plongement complexe $\rho:E\to\C$. On pose $\MC(G)=\MC(G,h)$.

\paragraph{Définition 3.1.} Un modèle sur $E$ de $\MC(G)$ consiste en
\begin{enumerate}[label=(\alph*),leftmargin=2.4em]
\item un schéma $M$ sur $E$, muni d'une action continue de $G(\Af)$;
\item un isomorphisme $m$ de $M\otimes_{E,\rho}\C$ avec $\MC(G)$ compatible à l'action de $G(\Af)$.
\end{enumerate}
Soit $F$ une extension finie de $E$, munie d'un plongement complexe qui prolonge celui de $E$. Si $M_E(G,h)$ est un modèle de $\MC(G,h)$ sur $E$, on désigne par $M_F(G,h)$ le modèle $M_E(G,h)\otimes_E F$ de $\MC(G,h)$ sur $F$.

\paragraph{Remarque 3.2.} Se donner un schéma $M$ sur $E$, muni d'une action continue de $G(\Af)$, revient à se donner
\begin{enumerate}[label=(\alph*),leftmargin=2.4em]
\item pour tout sous-groupe compact ouvert $K$ de $G(\Af)$, un schéma ${}_K M$ sur $E$;
\item pour $K$ et $L$ deux sous-groupes compacts ouverts de $G(\Af)$ et pour $x\in G(\Af)$ tel que $xKx^{-1}\subset L$, un homomorphisme
\[
J_{L,K}(x):{}_K M\longrightarrow {}_L M,
\]
ces homomorphismes vérifiant
\begin{align*}
(\alpha)\quad&J_{M,L}(y)J_{L,K}(x)=J_{M,K}(yx),\tag{3.2$\alpha$}\\
(\beta)\quad&J_{K,K}(x)=\id\qquad\text{si }x\in K,\tag{3.2$\beta$}\\
(\gamma)\quad&\text{pour }K\text{ distingué dans }L,\; J\text{ définit une action de }L/K\text{ sur }{}_K M,\\
&\text{et }J_{L,K}(e)\text{ définit }(L/K)\backslash {}_K M\xrightarrow{\sim}{}_L M.\tag{3.2$\gamma$}
\end{align*}
\end{enumerate}
On aura
\[
M=\limproj_K {}_K M,\tag{3.2.1}
\]
et
\[
{}_K M=K\backslash M.\tag{3.2.2}
\]

\paragraph{3.3.} Soit $M(G)$ un modèle sur $E$ de $\MC(G)$. Pour $K$ compact ouvert dans $G(\Af)$, on pose ${}_K M(G)=K\backslash M(G)$. Les ${}_K M(G)$ sont des schémas quasi-projectifs sur $E$, non nécessairement géométriquement connexes. L'isomorphisme structural $m$ induit des isomorphismes
\[
{}_K M(G)\otimes_{E,\rho}\C\xrightarrow{\sim}{}_K\MC(G).\tag{3.3.1}
\]

\paragraph{3.4.} Supposons que $G$ vérifie les hypothèses de 2.1, et soit $(M,m)$ un modèle de $\MC(G,h)$ sur $E$. Soit $\overline E$ la clôture algébrique de $E$ dans $\C$. Puisque $M$ est défini sur $E$, le groupe de Galois $\Gal(\overline E/E)$ agit sur l'ensemble profini
\[
\pi_0(M\otimes_E\overline E)=\limproj_K\pi_0({}_K M\otimes_E\overline E)\xrightarrow{\sim}_{m}\pi_0(\MC(G,h)).
\]
Le groupe $G(\Af)$ agit sur $\pi_0(M\otimes_E\overline E)$ via son quotient $\pi(G)/\pi_0(K_\infty)$, et $\pi_0(M\otimes_E\overline E)$ est un espace principal homogène sous le groupe commutatif $\pi(G)/\pi_0(K_\infty)$ (2.6). L'action de $\Gal(\overline E/E)$ commute à celle de $G(\Af)$, donc à celle de $\pi(G)/\pi_0(K_\infty)$. L'action d'un élément $\sigma$ de $\Gal(\overline E/E)$ ne peut donc être que la translation définie par un élément $\lambda(\sigma)$ de $\pi(G)/\pi_0(K_\infty)$, et $\lambda$ est un homomorphisme
\[
\lambda_M:\Gal(\overline E/E)\longrightarrow \pi(G)/\pi_0(K_\infty).\tag{3.4.1}
\]
Supposons que $E$ soit un corps de nombres. La théorie du corps de classe identifie alors le plus grand quotient abélien de $\Gal(\overline E/E)$ à $\pi_0(E^*(\A)/E^*(\Q))$, et (3.4.1) à
\[
\lambda_M:\Gal(\overline E/E)^{\ab}=\pi_0(E^*(\A)/E^*(\Q))\longrightarrow\pi(G)/\pi_0(K_\infty).\tag{3.4.2}
\]
Dans le cas particulier où $G'$ est simplement connexe, l'homomorphisme (3.4.2) s'identifie (2.7.1) à
\[
\lambda_M:\pi_0(E^*(\A)/E^*(\Q))\longrightarrow\pi_0(T(\A)/T(\Q))/\pi_0(K_\infty).\tag{3.4.3}
\]

\paragraph{Définition 3.5.} Les homomorphismes (3.4.1), (3.4.2), (3.4.3) s'appellent la loi de réciprocité du modèle $M$. Si $G'$ est simplement connexe, et que $\lambda_M$ provient d'un homomorphisme de groupes algébriques sur $\Q$ de $E^*$ dans $T$, ce dernier s'appellera encore la «loi de réciprocité» de $M$.

\paragraph{3.6.} Il existe une règle, applicable à chacun des modèles construits par Shimura, pour déterminer $E$ et la loi de réciprocité $\lambda_M$ à partir de $G$ et $h$. Le corps $E$, pour $G$ quelconque, et $\lambda_M$, pour $G'$ simplement connexe, se décrivent comme suit.

\paragraph{3.7.} L'homomorphisme composé $hr$ (1.3)
\[
hr:\mathbb G_{m,\C}\xrightarrow{r}\bS_\C\xrightarrow{h}G_\C
\]
est un homomorphisme sur $\C$ entre groupes algébriques définis sur $\Q$. Le sous-corps $E$ de $\C$ est le corps de définition $E(G,h)$ de la classe de conjugaison de cet homomorphisme.

Pour $G$ semi-simple adjoint, le corps $E(G,h)$ peut se décrire comme suit. Soient $\Delta$ le diagramme de Dynkin de $G_\C$ et $|\Delta|$ l'ensemble des sommets de $\Delta$. Le groupe de Galois $\Gal(\C/\Q)$ agit sur $\Delta$.

Soit $H$ un tore maximal de $G_\C$, muni d'un système de racines simples $(\alpha_i)_{i\in |\Delta|}$. Les $\alpha_i$ identifient $H$ à $\Gm^{\Delta}$. Pour tout homomorphisme $u:\mathbb G_{m,\C}\to G_\C$, il existe alors un unique homomorphisme $u':\mathbb G_{m,\C}\to H=\mathbb G_{m,\C}^{\Delta}$,
\[
x\longmapsto (x^{n_i})_{i\in |\Delta|}\qquad(n_i\ge0),
\]
qui soit conjugué à $u$. La construction $u\mapsto \eta(u)=(n_i)_{i\in |\Delta|}$ identifie l'ensemble des classes de conjugaison d'applications de $\mathbb G_{m,\C}$ dans $G_\C$ avec l'ensemble des fonctions de $|\Delta|$ dans $\mathbb N$. On en déduit que le corps $E(G,h)$ est défini par le sous-groupe de $\Gal(\C/\Q)$ qui stabilise $\eta(hr)$.

\paragraph{Proposition 3.8.} Soit
\[
G/C=\prod_{i=1}^g G_i
\]
la décomposition du groupe adjoint $G/C$ en facteurs $\Q$-simples, et soit $h_i$ le composé
\[
h_i:\bS\xrightarrow{h}G_\R\longrightarrow (G_i)_\R.
\]
(i) Le sous-corps $E(G,h)$ de $\C$ est le composé des sous-corps $E(T,\nu h)$ et $E(G_i,h_i)$ ($1\le i\le g$) de $\C$.

(ii) Si l'involution d'opposition du diagramme de Dynkin de $G_i$ respecte $\eta(h_ir)$, alors $E(G_i,h_i)$ est totalement réel. Sinon, $E(G_i,h_i)$ est une extension quadratique totalement imaginaire d'un corps de nombres totalement réel.

L'assertion (i) est facile à vérifier, et il suffit de prouver (ii) pour $G$ $\Q$-simple et adjoint. Soit $\Delta$ son diagramme de Dynkin. Il résulte des hypothèses 1.5 que $G_\R$ admet un tore maximal compact; ceci implique que la conjugaison complexe agit sur $\Delta$ par l'involution d'opposition $\iota$. Celle-ci est dans le centre du groupe des automorphismes de $\Delta$.

Soit $E'$ l'extension galoisienne de $\Q$ définie par le sous-groupe de $\Gal(\C/\Q)$ qui agit trivialement sur $\Delta$. Alors, l'image $c$ de la conjugaison complexe est centrale dans $\Gal(E'/\Q)$; il en résulte que $E'$ est totalement réel ou extension quadratique totalement imaginaire d'un corps totalement réel. Enfin, $c$ est l'identité sur $E(G,h)\subset E'$ si et seulement si $\iota$ respecte $\eta(hr)$, d'où l'assertion.

\paragraph{3.9.} On reprend les notations de 3.6. On suppose que $G'$ est simplement connexe, que $E$ contient $E(G,h)$, et que le plongement complexe de $E$ induit celui de $E(G,h)$. Le morphisme composé
\[
r''(h):\mathbb G_{m,\C}\xrightarrow{r}\bS_\C\xrightarrow{h}G_\C\xrightarrow{\nu}T_\C
\]
ne dépend que de la classe de conjugaison de $hr$. Il est donc défini sur $E$, i.e. provient par extension des scalaires de $E$ à $\C$ d'un homomorphisme de groupes algébriques sur $E$, encore noté
\[
r''(h):\Gm\longrightarrow T_E.
\]
Appliquons la restriction des scalaires à la Weil, et soit $r'(h):E^*\to T$ le composé
\[
E^*=\Res_{E/\Q}\Gm\xrightarrow{\Res_{E/\Q}(r''(h))}\Res_{E/\Q}(T_E)\xrightarrow{N_{E/\Q}}T.
\]
La loi de réciprocité de $M$ est l'inverse $r(G,h)$ de $r'(h)$ :
\[
r(G,h)=r'(h)^{-1}:E^*\longrightarrow T .\tag{3.9.1}
\]

\paragraph{3.10. Exemple II.} Soient $H$ un tore et $h:\bS\to H_\R$. Les conditions de 1.5 sont automatiquement vérifiées. Les ${}_K\MC(H,h)$ sont des ensembles finis. Un schéma réduit fini sur le corps $E(H,h)$ (3.7) s'identifie à un ensemble galoisien. On en déduit qu'il existe, à isomorphisme unique près, un et un seul modèle de $\MC(H,h)$ sur $E(H,h)$, de loi de réciprocité donnée par 3.9.1. On le note $M(H,h)$.

\paragraph{3.11.} Soit $F$ une extension finie de $E$, munie d'un plongement complexe qui prolonge celui de $E$. Le diagramme
\[
\begin{array}{ccc}
F^*&\xrightarrow{N_{F/E}}&E^*\\
&\searrow&\downarrow\\
&&T
\end{array}
\]
est commutatif. Si $M_E(G,h)$ est un modèle sur $E$ de $\MC(G,h)$, ayant (3.9.1) pour loi de réciprocité, il en résulte que $M_F(G,h)$ (3.1) a encore (3.9.1) pour loi de réciprocité.

\paragraph{3.12.} Soient $u:(G_1,h_1)\to(G_2,h_2)$ comme en 1.14, et $M^i(G_i,h_i)$ un modèle de $\MC(G_i,h_i)$ sur $E_i$. Soit $E$ le composé des sous-corps $E_1$ et $E_2$ de $\C$. On a (notation de 3.1)
\[
M^i(G_i,h_i)\otimes_{E_i}E=M_E(G_i,h_i)\qquad(i=1,2).
\]
Il a donc un sens de se demander si
\[
u:\MC(G_1,h_1)\longrightarrow\MC(G_2,h_2)
\]
est défini sur $E$.

\paragraph{Définition 3.13.} Soient $G$ un groupe réductif connexe sur $\Q$ vérifiant la condition de 2.1, $h:\bS\to G_\R$ un homomorphisme vérifiant les conditions de 1.5, et $E$ une extension de $E(G,h)$ munie d'un plongement complexe prolongeant celui de $E(G,h)$. Un modèle $M_E(G,h)$ de $\MC(G,h)$ sur $E$ est dit faiblement canonique si, pour tout tore $u:H\to G$ dans $G$, muni de $h':\bS\to H_\R$ tel que $uh'$ soit le conjugué de $h$ par un élément de $G(\R)$, le morphisme (1.14)
\[
\MC(H,h')\longrightarrow\MC(G,h)
\]
est défini (3.10), (3.12) sur le composé $E(H,h')\cdot E\subset\C$. Un modèle $M_E(G,h)$ est dit canonique s'il est faiblement canonique et que $E=E(G,h)$.

\paragraph{3.14. Traduction.} Soit $M$ un modèle canonique de $\MC(G,h)$, de loi de réciprocité $\lambda_M$. Si $G'$ est simplement connexe, alors $\lambda_M$ est donnée par 3.9.1 (5.6). Pour tout sous-groupe compact ouvert $K$ de $G(\Af)$, désignons encore par $\nu(K)$ l'image de $K$ dans $\pi(G)/\pi_0(K_\infty)$, et soit $E(K)$ l'extension de $E(G,h)$ définie par le groupe de classes d'idèles $\lambda_M^{-1}(\nu(K))$. On vérifie que $E(K)$ est le corps de définition de la composante neutre ${}_K\MC^0$ de ${}_K\MC$, considérée comme sous-schéma de ${}_K\MC={}_K M\otimes_E\C$.

Par définition, ${}_K\MC^0$ provient par extension des scalaires de $E(K)$ à $\C$ d'un sous-schéma ${}_K M'$ de ${}_K M\otimes_EE(K)$. Le schéma ${}_K M'/E(K)$ est géométriquement connexe sur $E(K)$, et on peut encore décrire le schéma ${}_K M'$ comme celle des composantes connexes (sur $E$) de ${}_K M$ qui, après extension des scalaires de $E$ à $\C$, contient l'origine
\[
\begin{array}{ccc}
{}_K\MC^0&\subset&{}_K\MC\\
\downarrow&&\downarrow\\
\Spec E(K)&\longrightarrow&\Spec\C\\
\downarrow&&\\
\Spec E&=&\Spec E.
\end{array}
\]
Shimura a l'habitude de s'exprimer en termes du système des variétés ${}_K M'/E(K)$ et des homomorphismes de type 2.2 qu'on a entre celles-ci. Pour un énoncé typique, voir [14, p. 146]. Pour la traduction, cf. [14, 2.7].




\paragraph{3.15.} Pour $u:(H,h')\longrightarrow(G,h)$ comme en 3.13, l'image de $M_{\C}(H,h')$ dans ${}_K M_{\C}(G,h)$ est un ensemble fini. L'hypothèse est que cette partie de ${}_K M_{\C}(G,h)$ soit définie sur $E(H,h')$, que ses points soient définis sur une extension abélienne de $E(H,h')$, et que $\Gal(\overline{E(H,h')}/E(H,h'))$ les permute d'une manière prescrite.

On appellera \emph{spéciaux} les points de ${}_K M_{\C}(G,h)$ dans l'image d'un $M_{\C}(H,h')$.

Supposons connu un plongement projectif
\[
q:{}_K M_{\C}^{0}\longrightarrow \mathbb P^r,
\]
qui soit défini sur $E(K)$. On peut l'interpréter comme
\[
\widetilde q:X^0\longrightarrow \mathbb P^r
\]
(cf. 1.7, 1.8). Si $uh'\in X$ est dans $X^0$, alors les coordonnées de $\widetilde q(uh')$ engendrent sur $E(K)E(H,h')$ une extension abélienne de $E(H,h')$ qu'on peut expliciter comme corps de classe.

\paragraph{3.16.} L'exemple classique d'une telle situation est fourni par le schéma de module des courbes elliptiques, correspondant à $G=\GL_2$, $K=\GL_2(\widehat{\Z})$, $h$ comme en 1.6, et par ceux de ses points définis par des courbes elliptiques à multiplication complexe.

Soit $j$ l'invariant modulaire, vu comme fonction sur le demi-plan de Poincaré $X^0$. Si $\tau\in X^0$ appartient à un corps quadratique imaginaire $K$, soit
\[
L(\tau)=\Z\oplus \Z\tau\subset\C,
\qquad
\mathcal O(\tau)=\{\,k\in K\mid kL(\tau)\subset L(\tau)\,\}.
\]
Le corps $K(j(\tau))$ est l'extension abélienne de $K$ de groupe de Galois le groupe des classes d'idéaux de $\mathcal O(\tau)$. Si $\sigma$ dans Galois correspond à un module inversible $L$, et si
\[
L(\tau')\simeq L(\tau)\otimes_{\mathcal O(\tau)}L,
\]
on a $j(\tau)=j(\tau')^\sigma$.

Pour des exemples explicites plus étranges, voir [15], p. 45--46.

\section*{\S 4. Variétés abéliennes}

\paragraph{4.1.} Soit $k$ un corps algébriquement clos de caractéristique $0$. Soit $A$ une variété abélienne sur $k$. Pour $n\in\mathbb N^+$, on désigne par $A_n$ le noyau de la multiplication par $n$. Les $A_n$ forment un système projectif
\[
\varphi_{nm,n}:A_{nm}\longrightarrow A_n,\qquad x\longmapsto x^m,
\]
et on pose
\[
\widehat T(A)=\varprojlim_n A_n,\qquad
\widehat V(A)=\Af\otimes_{\widehat{\Z}}\widehat T(A).
\]
Le $\widehat{\Z}$-module $\widehat T(A)$ est le produit des modules de Tate, ou représentations de Weil, $T_\ell(A)$. Si $k=\C$, on a
\[
\tag{4.1.1}
\widehat T(A)=\widehat{\Z}\otimes H_1(A,\Z),
\]
\[
\tag{4.1.2}
\widehat V(A)=\Af\otimes H_1(A,\Q).
\]

\paragraph{4.2.} La catégorie des \emph{variétés abéliennes à isogénie près} sur $k$ est la catégorie dont les objets sont les variétés abéliennes sur $k$ et dans laquelle
\[
\Hom_{\mathrm{iso}}(A,B)=\Hom(A,B)\otimes\Q.
\]
On désigne par $A\otimes\Q$ la variété abélienne à isogénie près ``sous-jacente'' à une variété abélienne $A$. Tout foncteur additif de la catégorie des variétés abéliennes dans une catégorie additive $\Q$-linéaire se factorise par le foncteur $A\mapsto A\otimes\Q$. Ainsi, le foncteur contravariant ``variété duale'' $A\mapsto A^\vee$ se prolonge aux variétés abéliennes à isogénie près. Le foncteur $A\mapsto\widehat V(A)$ se prolonge de même. De plus, si $A_0$ est une variété abélienne à isogénie près, il revient au même de se donner $B$, plus un isomorphisme $B\otimes\Q=A_0$, ou de se donner $\widehat T(B)\subset\widehat V(A_0)$.

\paragraph{4.3.} Soit $A$ une variété abélienne sur $k$. On désigne par $\NS(A)$ le groupe des classes d'équivalence algébrique de faisceaux inversibles sur $A$. On définit une \emph{polarisation effective}, respectivement une \emph{polarisation}, de $A$ comme un élément de $\NS(A)$, respectivement de $\NS(A)\otimes\Q$, dont un multiple positif est défini par un plongement projectif de $A$. Une \emph{polarisation homogène} est un élément de $(\NS(A)\otimes\Q)/\Q^*$ qui est la classe d'une polarisation.

Rappelons que $\NS(A)\otimes\Q$ s'identifie, par une bijection $p\mapsto p'$, à l'ensemble des éléments de $\Hom(A,A^\vee)\otimes\Q$ égaux à leur transposé. Un isomorphisme
\[
u\in\Hom(A\otimes\Q,B\otimes\Q)
\]
induit des isomorphismes de $\Hom(A,A^\vee)\otimes\Q$ avec $\Hom(B,B^\vee)\otimes\Q$, de $\NS(A)\otimes\Q$ avec $\NS(B)\otimes\Q$, et transforme polarisation en polarisation. Ceci permet de parler de polarisation d'une variété abélienne à isogénie près.

\paragraph{4.4.} Une polarisation $p$ de $A$ définit dans $\End(A)\otimes\Q$ une involution positive
\[
u\longmapsto p'^{-1}u^*p',
\]
qui ne dépend que de la polarisation homogène $\Q^*p$.

Soit $F$ un produit de corps totalement réels, et soit $\rho:F\longrightarrow\End(A)\otimes\Q$. Soit $\NS_\rho(A)\subset\NS(A)$ l'ensemble des $p$ tels que
\[
p'\rho(f)=\rho(f)^*p'\qquad(f\in F).
\]
Le vectoriel $\NS_\rho(A)\otimes\Q$ est muni, pour $f\in F$, d'une action de $F$:
\[
(f.p)'=p'\rho(f)=\rho(f)^*p'.
\]
Une \emph{polarisation faible} de $A$ relative à $\rho,F$ est un élément de
\[
(\NS_\rho(A)\otimes\Q)/F^*
\]
qui est la classe d'une polarisation. L'ensemble des polarisations faibles de $A$ ne dépend que de $A\otimes\Q$. La restriction au centralisateur de $F$ de l'involution de $\End(A)\otimes\Q$, définie par une polarisation $p\in\NS_\rho(A)$, ne dépend que de la polarisation faible $F^*p$.

\paragraph{4.5.} Si $A_0^\vee$ est la duale d'une variété abélienne à isogénie près $A_0$, alors $\widehat V(A_0)$ et $\widehat V(A_0^\vee)$ sont en dualité à valeurs dans $\Af(1)$. Si $p$ est une polarisation de $A_0$, la \emph{forme de polarisation}
\[
\psi_p(x,y)=\langle x,p'(y)\rangle
\]
est une forme alternée non dégénérée sur le $\Af$-module libre $\widehat V(A_0)$, à valeurs dans $\Af(1)$.

Les rappels 4.1 et 4.5 s'étendent aux schémas abéliens sur une base quelconque.

\paragraph{4.6.} Faisons $k=\C$. Le foncteur additif $A\mapsto H_1(A,\Q)$ se prolonge aux variétés abéliennes à isogénie près sur $\C$ (4.2). Pour $A_0$ une variété abélienne à isogénie près, $H_1(A_0,\Q)$ et $H_1(A_0^\vee,\Q)$ sont en dualité. Via 4.1.2 et l'isomorphisme de $\Af(1)$ avec $\Af$ défini par l'exponentielle, cette dualité est compatible à celle considérée en 4.5. La forme alternée
\[
\psi_p(x,y)=\langle x,p'(y)\rangle
\]
sur $H_1(A_0,\Q)$ définie par une polarisation $p$ de $A_0$ s'appelle encore la forme de polarisation.

La structure de Hodge de $H_1(A_0,\Q)$ est l'homomorphisme
\[
h:\bS\longrightarrow \GL(H_1(A_0,\Q))_\R,
\]
tel que $\bS$ agisse sur $H_1(A_0,\Q)\otimes\C$ par les caractères $z^{-1}$ et $\bar z^{-1}$, et que $F(h)^\circ$ (1.5) soit le noyau de l'exponentielle
\[
H_1(A_0,\Q)\otimes\C\longrightarrow\Lie(A).
\]

\paragraph{Théorème 4.7 (Riemann).} Les constructions
\[
(A,p)\longmapsto (H_1(A,\Q),\psi_p,H_1(A,\Z),h),
\]
\[
(V,\psi,V_\Z,h)\longmapsto
F(h)^\circ\backslash (V\otimes\C)/V_\Z
\]
établissent une équivalence entre
\begin{enumerate}[label=(\alph*)]
\item variétés abéliennes polarisées sur $\C$;
\item systèmes formés d'un vectoriel $V$ sur $\Q$, d'une forme alternée non dégénérée $\psi$ sur $V$, d'un réseau entier $V_\Z$ dans $V$ et d'un homomorphisme $h$ de $\bS$ dans le groupe $\Gp(V)_\R$ des similitudes symplectiques de $V\otimes\R$, du type 1.6 et tel que
\[
\psi(x,h(i)x)>0\qquad(x\in V\otimes\R,\ x\ne0).
\]
\end{enumerate}

\paragraph{Variante 4.8.} Il revient au même de se donner une variété abélienne à isogénie près $A$ sur $\C$, munie d'une polarisation homogène, ou de se donner un vectoriel sur $\Q$, muni d'une forme alternée non dégénérée $\psi$, donnée à un facteur près, et de $h:\bS\longrightarrow\Gp(V)$ du type 1.6.

\paragraph{4.9.} Soient $L$ une algèbre semi-simple à involution sur $\Q$, et $V$ un vectoriel sur $\Q$, muni d'une structure de $L$-module fidèle et d'une forme alternée non dégénérée $\psi$ telle que
\[
\tag{4.9.1}
\psi(\ell x,y)=\psi(x,\ell^*y).
\]
On désigne par $G$ le groupe algébrique sur $\Q$ des similitudes symplectiques $L$-linéaires de $V$. Le groupe $G(\Q)$ est l'ensemble des $g\in\GL_L(V)$ tels qu'il existe $\mu(g)\in\Q^*$ vérifiant
\[
\tag{4.9.2}
\psi(gx,gy)=\mu(g)\psi(x,y).
\]
On suppose donné un homomorphisme $h_0:\bS\longrightarrow G_\R$ tel que, via $h_0$, $\bS$ agisse sur $V_\C$ par les caractères $z^{-1}$ et $\bar z^{-1}$ et que la forme $\psi(x,h_0(i)y)$ soit symétrique et définie positive. Les conditions (1.5.1) à (1.5.3) sont alors vérifiées par $(G,h_0)$, et l'involution de $L$ est positive.

\paragraph{4.10.} Appliquons (1.11) à $(G,h_0)$ et à la représentation $V$ de $G$. Une $G$-structure $\bar\xi$ sur un vectoriel $H$ s'interprète, par 1.8, comme la donnée sur $H$ d'une structure de $L$-module et d'une forme alternée $\psi$, donnée à un facteur près, le vectoriel $H$, muni de ces structures, étant isomorphe à $V$. Soit $\bar\xi$ une $G$-structure sur $H$. Une donnée 1.9 (b), $h:\bS\to G^1_\R$ sur $H$, définit alors, d'après 4.8 appliqué à $(H,\psi,h)$, une variété abélienne à isogénie près $A$, munie d'une polarisation homogène $\bar p$, telle que $H$ soit $H_1(A,\Q)$, $h$ la structure de Hodge de $H_1(A,\Q)$ et $\psi$ la forme de polarisation. De plus, la structure de $L$-module de $H$ provient de $\rho:L\to\End(A)$, et $H$, muni de $\bar\xi$ et $h$, est déterminé par $(A,\bar p,\rho)$.

Pour qu'un triple $(A,\bar p,\rho)$ provienne d'un $H$ comme plus haut, il faut et il suffit que
\[
\tag{4.10.1}
H_1(A,\Q),\text{ muni de sa structure de }L\text{-module et de la forme de polarisation donnée à un facteur près, soit isomorphe à }V;
\]
\[
\tag{4.10.2}
\text{pour }t:V\longrightarrow H_1(A,\Q)\text{ un isomorphisme, }t^{-1}h_1t\text{ soit conjugué à }h_0.
\]
Soit $t$ l'application de $L$ dans $\C$ donnée par
\[
\tag{4.10.3}
t(\ell)=\Tr(\ell;V_\C/F^0(h_0)).
\]

\paragraph{Scholie 4.11.} Soit $K$ un sous-groupe compact ouvert de $G(\Af)$. Les points de ${}_K M_{\C}(G,h_0)$ correspondent bijectivement aux classes d'isomorphie de variétés abéliennes à isogénie près $A$, munies de
\begin{enumerate}[label=(\alph*)]
\item $\rho:L\to\End(A)$ tel que
\[
\Tr(\rho(\ell),\Lie(A))=t(\ell)\qquad(\ell\in L),
\]
\item une polarisation homogène $\bar p$ qui induise l'involution donnée de $L$,
\item une classe modulo $K$, $\bar k$, de similitudes symplectiques $L$-linéaires
\[
k:\widehat V(A)\xrightarrow{\sim}V\otimes\Af,
\]
\end{enumerate}
et telles que
\[
(*)\qquad\text{les conditions (4.10.1) et (4.10.2) soient vérifiées.}
\]
On applique 1.11 et 4.10, en notant qu'une donnée 1.9 (c) s'identifie à une donnée (c), et que les conditions en (b) et (a) résultent de (4.10.1) et (4.10.2) respectivement.

\paragraph{4.12.} Soient $K$ un sous-groupe compact ouvert de $G(\Af)$, $V_\Z$ un réseau entier dans $V$ tel que $V_{\widehat{\Z}}=V_\Z\otimes\widehat{\Z}$ soit $K$-invariant, $L_\Z$ un ordre de $L$ tel que $L_\Z V_\Z\subset V_\Z$, $\psi_\Z$ un multiple rationnel positif de la forme alternée de $V$, à valeurs entières sur $V_\Z$, et $V'_\Z$ le plus grand réseau de $V$ tel que $\psi_\Z(V_\Z,V'_\Z)\subset \Z(1)$. Soient $n$ un entier, et
\[
K_n=\{\,g\in G(\Af)\mid (g-1)V'_{\widehat{\Z}}\subset nV_{\widehat{\Z}}\,\}.
\]
On suppose $n$ assez grand pour que $K\supset K_n$.

Si $A$ est comme en 4.11, alors $k^{-1}(V_{\widehat{\Z}})\subset\widehat V(A)$ ne dépend pas de $k\in\bar k$, et définit une variété abélienne $B$ avec $B\otimes\Q\simeq A$ et
\[
\widehat T(B)=k^{-1}(V_{\widehat{\Z}})\subset\widehat V(A).
\]
L'action de $L$ sur $A$ induit une action de $L_\Z$ sur $B$. Il existe une unique polarisation effective $p$ de $B$, avec $p\in\bar p$, telle que $k^{-1}(V'_{\widehat{\Z}})$ soit le plus grand ``réseau'' $\widehat T'(B)\supset\widehat T(B)$ vérifiant
\[
\psi_p(\widehat T(B),\widehat T'(B))\subset \widehat{\Z}(1).
\]
Enfin, l'ensemble $\bar k$ d'isomorphismes symplectiques $L$-linéaires de $\widehat T(B)$ avec $V_{\widehat{\Z}}$ est l'image réciproque de son image $\bar k_n$ dans $\Isom(B_n,V_\Z/nV_\Z)$.

Ceci permet encore d'interpréter les points de ${}_K M_{\C}(G,h_0)$ comme correspondant aux classes d'isomorphie de systèmes $(B,p,\rho,\bar k_n)$ consistant en
\begin{enumerate}[label=(\alph*)]
\item une variété abélienne polarisée $(B,p)$, à multiplication complexe par $L_\Z$, vérifiant 4.11 (a) (b),
\item une classe modulo $K/K_n$, $\bar k_n$, d'isomorphismes
\[
k_n:B_n\xrightarrow{\sim}V_\Z/nV_\Z,
\]
qui peuvent se relever en des isomorphismes symplectiques $L$-linéaires
\[
k:\widehat T(B)\longrightarrow V_{\widehat{\Z}},
\]
et qui vérifient (4.10.1) (4.10.2).
\end{enumerate}

\paragraph{4.13.} Désignons par $F$ l'algèbre des invariants par $*$ du centre de $L$. C'est un produit de corps totalement réels. Il existe une et une seule forme $F$-bilinéaire alternée $\psi_F$ sur $V$ telle que
\[
\Tr_{F/\Q}\psi_F=\psi.
\]
On désigne par $G_1$ le groupe des similitudes symplectiques $L$-linéaires de $\psi_F$: $G_1(\Q)$ est l'ensemble des $g\in\GL_L(V)$ tels qu'il existe $\mu(g)\in F^*$ vérifiant
\[
\tag{4.13.1}
\psi_F(gx,gy)=\mu(g)\psi_F(x,y),
\]
c'est-à-dire
\[
\tag{4.13.2}
\psi(gx,gy)=\psi(\mu(g)x,y).
\]
Comme plus haut, on vérifie la variante suivante.

\paragraph{Variante 4.14.} Soit $K$ un sous-groupe compact ouvert de $G_1(\Af)$. Les points de ${}_K M_{\C}(G_1,h_0)$ correspondent bijectivement aux classes d'isomorphie de variétés abéliennes à isogénie près $A$, munies de
\begin{enumerate}[label=(\alph*)]
\item $\rho:L\longrightarrow\End(A)$ comme en 4.11,
\item une polarisation faible, relative à $F$, $\bar p$, qui induise l'involution donnée de $L$,
\item une classe modulo $K$ de similitudes symplectiques $L$-linéaires
\[
k:\widehat V(A)\xrightarrow{\sim}V\otimes\Af
\]
pour les formes $F\otimes\Af$-bilinéaires des deux membres,
\end{enumerate}
et telles que
\[
(*)\qquad\text{des conditions analogues à (4.10.1) et (4.10.2) soient vérifiées.}
\]

\paragraph{4.15.} On peut préciser 4.11 et 4.14 en montrant que ${}_K M_{\C}(G,h)$ est le schéma de modules grossier d'un foncteur évident sur les schémas de type fini sur $\C$.

\paragraph{4.16. Exemple I} (suite de 1.6, 1.11). Plaçons-nous dans le cas particulier de 4.9 où $L=\Q$ et où $\dim(V)=2g$. On a alors $G=\Gp(V)$, le groupe des similitudes symplectiques. Pour $h_0:\bS\to G_\R$ de type 1.6, on a $E(\Gp,h_0)=\Q$. Soient $V_\Z$, $n$ et $K_n$ comme 1.11. Pour tout schéma $S$, soit $F(S)$ l'ensemble des classes d'isomorphie de schémas abéliens polarisés de la série principale $A/S$, munis d'une similitude symplectique
\[
k_n:A_n\xrightarrow{\sim}(V_\Z/nV_\Z)_S.
\]
Pour $n\ge3$, les objets classifiés n'ont plus d'automorphismes et le foncteur $F$ est représenté par un $\Q$-schéma ${}_{K_n}M$ [6]. D'après 1.11, 4.11 et 4.12, on dispose de
\[
m_n:{}_{K_n}M(\C)\xrightarrow{\sim}{}_{K_n}M_{\C}(\Gp,h_0).
\]
Dans le langage 3.1, on a

\paragraph{Proposition 4.17.} Le $\Q$-schéma
\[
M(\Gp,h_0)=\varprojlim_{K_n}{}_{K_n}M
\]
est un modèle sur $\Q=E(\Gp,h_0)$ de $M_{\C}(\Gp,h_0)$.

\paragraph{4.18.} Plaçons-nous dans le cas particulier de 4.9 où $V$ est un $L$-module monogène. Les groupes $G$ et $G_1$ sont alors contenus dans le centre de $L^*$.

Soient $\bar E$ une clôture algébrique de $E(G,h)$ et $M^+$ l'ensemble des classes d'isomorphie $[A,\rho,k,\bar p]$ d'objets $(A,\rho,k,\bar p)$ consistant en
\begin{enumerate}[label=(\alph*)]
\item une variété abélienne à isogénie près $A/\bar E$, munie d'une polarisation homogène $\bar p$ et de $\rho:L\to\End(A)$ tel que, avec la notation de 4.10.3, on ait pour $\ell\in L$
\[
\Tr(\rho(\ell),\Lie(A))=t(\ell)\in E(G,h);
\]
\item un isomorphisme $L\otimes\Af$-linéaire
\[
k:\widehat V(A)\xrightarrow{\sim}V\otimes\Af.
\]
\end{enumerate}
La variété abélienne $A$ est donc ``de type CM''. On renvoie à [16] pour la démonstration du théorème fondamental suivant.

\paragraph{Théorème 4.19 (Shimura--Taniyama).} Le groupe de Galois $\Gal(\bar E/E(G,h))$ agit sur $M^+$ via son plus grand quotient abélien
\[
\pi_0\bigl(E(G,h)^*(\A)/E(G,h)^*(\Q)\bigr).
\]
Pour $e\in E(G,h)^*(\A)$, d'image $\varphi(e)$ dans Galois rendu abélien, et de composante $e^f$ dans $E(G,h)^*(\Af)$, on a
\[
\tag{4.19.1}
\varphi(e)([A,\rho,k,\bar p])=[A,\rho,r(G,h)(e^f)k,\bar p].
\]

Soient $\tau$ un plongement de $\bar E$ dans $\C$ qui prolonge celui de $E(G,h)$, $K$ un sous-groupe compact ouvert de $G(\Af)$ et ${}_K M(G,h)(\bar E)$ l'ensemble des classes d'isomorphie d'objets $(A,\rho,\bar k,\bar p)$ où
\begin{enumerate}[label=(\alph*)]
\item $A,\rho$ et $\bar p$ sont comme plus haut;
\item $\bar k$ est une classe modulo $K$ de similitudes symplectiques $L\otimes\Af$-linéaires
\[
k:\widehat V(A)\xrightarrow{\sim}V\otimes\Af;
\]
\item le vectoriel $H_1(A_\C,\Q)$, où $A_\C$ est défini via $\tau$, muni de l'action de $L$ et de la forme de polarisation donnée à un facteur près, est isomorphe à $V$.
\end{enumerate}
D'après 4.19, la condition (c) est indépendante du choix de $\tau$, de sorte que l'ensemble fini ${}_K M(G,h)(\bar E)$ est muni d'une action de $\Gal(\bar E/E(G,h))$; en tant qu'ensemble galoisien, il s'identifie à l'ensemble des points sur $\bar E$ d'un $E(G,h)$-schéma fini étale ${}_K M(G,h)$. D'après 4.14, ${}_K M(G,h)(\C)=G(\R)XK\backslash G(\A)/G(\Q)$, et par 4.19 on a:

\paragraph{Proposition 4.20.} Le $E(G,h)$-schéma
\[
M(G,h)=\varprojlim_K {}_K M(G,h)
\]
est un modèle canonique de $M_{\C}(G,h)$.

\paragraph{Théorème 4.21.} Le modèle $M(\Gp,h_0)$ construit en 4.17 est un modèle canonique.

Soient $(V,\psi)$ comme en 4.16, $L$, $\rho:L\to\End(V)$, $G$ et $h:\bS\to G_\R$ comme en 4.18. On dispose de
\[
u:(G,h)\longrightarrow(\Gp,h_0).
\]
Soit $K\subset G(\Af)$ tel que $u(K)\subset K_n$ (4.16). Si $\bar E$ est une clôture algébrique de $E(G,h)$, on définit
\[
u:{}_K M(G,h)(\bar E)\longrightarrow{}_{K_n}M(\Gp,h_0)(\bar E)
\]
par ``oubli de la multiplication complexe'', en associant à $[A,\rho,\bar k,\bar p]$ le triple formé de
\begin{enumerate}[label=(\alph*)]
\item la variété abélienne $B$, munie de $B\otimes\Q\simeq A$ et telle que
\[
\widehat T(B)=k^{-1}(V_{\widehat{\Z}})\subset\widehat V(A)\qquad(k\in\bar k);
\]
\item la polarisation $\bar p$ de $B$;
\item l'isomorphisme
\[
k:B_n\xrightarrow{\sim}V_\Z/nV_\Z\qquad(k\in\bar k).
\]
\end{enumerate}
Les applications $u$ définissent un morphisme de modèles
\[
u:M(G,h)\longrightarrow M(\Gp,h_0)\otimes E(G,h).
\]
On achève la démonstration en montrant ([8]) que, pour tout
\[
u:(H,h')\longrightarrow(\Gp,h_0)
\]
comme en 3.13, il existe $v:(G,h)\longrightarrow(\Gp,h_0)$ comme plus haut, pour $L$ convenable, tel que $uh'=vh$.

\paragraph{Variante 4.22.} Soient $F$ un corps de nombre totalement réel, $n=[F:\Q]$, et $\Gp_{2g}^{(F)}$ le groupe des similitudes symplectiques d'un $F$-vectoriel symplectique de dimension $2g$. Soit
\[
h_0:\bS\longrightarrow \Gp_{2g,\R}^{(F)}\simeq \Gp_{2g}(\R)^n
\]
comme en 1.6. On peut construire un modèle canonique de $M_{\C}(\Gp_{2g}^{(F)},h_0)$ à partir de modules des schémas abéliens faiblement polarisés.





\section*{\S 5. Techniques de construction}

La définition 3.13 n'est utilisable que parce qu'il y a « beaucoup » de points spéciaux.

\paragraph{Théorème 5.1.} Soient $G$ et $h$ comme en 3.13. Pour toute extension finie $F$ de $E(G,h)$, il existe $(H,h',u:H\hookrightarrow G)$ comme en 3.13 tels que l'extension $E(H,h')$ de $E(G,h)$ soit linéairement disjointe de $F$.

Soit
\[
Y=\Hom(\Gm,G)/G
\]
le schéma des classes de conjugaison de morphismes de $\Gm$ dans $G$. Le groupe de Galois $\Gal(\overline{\Q}/\Q)$ agit sur l'ensemble discret infini $Y(\C)$ des classes de conjugaison de morphismes de $\Gm{}_{\C}$ dans $G_{\C}$. Soit $Y_o$ le sous-schéma fini de $Y$ tel que $Y_o(\C)$ soit l'orbite sous Galois de la classe de $hr$. C'est le spectre de $E(G,h)$.

Soient $V'=\Lie(G)$, considéré comme espace affine sur $\Q$, et soit $V$ le sous-schéma ouvert des éléments réguliers de l'algèbre de Lie. Pour $v$ dans $V$, on désigne par $T_v$ le centralisateur de $v$, un tore maximal. Soit $W$ le schéma de module des tores maximaux $T$ de $G$, munis de
\[
s:\Gm\longrightarrow T,\qquad v\in\Lie(T),
\]
tels que
\begin{enumerate}[label=(\alph*)]
\item $v$ est régulier dans $\Lie(G)$;
\item la classe de $s:\Gm\to G$ est dans $Y_o$.
\end{enumerate}
Le morphisme
\[
f:W\longrightarrow V,\qquad (T,s,v)\longmapsto v
\]
est fini étale surjectif, noter que $T=T_v$. Soit $p$ le morphisme de $W$ dans $Y_o$,
\[
p:(T,s,v)\longmapsto \text{classe de }s.
\]
On obtient le diagramme
\[
\tag{5.1.1}
\begin{array}{ccc}
W & \xrightarrow{\ f\ } & V\\
{\scriptstyle p}\downarrow && \downarrow\\
\Spec(E(G,h)) & \longrightarrow & \Spec(\Q).
\end{array}
\]

\paragraph{Lemme 5.1.2.} Le schéma $W$ est irréductible, et les fibres géométriques de $p$ sont géométriquement irréductibles.

Puisque $Y_o$ est irréductible, le spectre d'une extension de $\Q$, il suffit de prouver la seconde assertion. Elle résulte de ce que
\begin{enumerate}[label=(\alph*)]
\item puisque $G$ est connexe, le schéma sur $\C$ des homomorphismes de $\Gm{}_{\C}$ dans $G_{\C}$ conjugués à un homomorphisme donné est irréductible;
\item pour $i:\Gm{}_{\C}\to G_{\C}$ donné, d'après le théorème de conjugaison des tores maximaux appliqué au centralisateur connexe de $i(\Gm{}_{\C})$, le schéma des tores maximaux de $G_{\C}$ contenant $i(\Gm{}_{\C})$ est irréductible.
\end{enumerate}

Soit $U\subset V(\R)$ l'ensemble des $v\in V(\R)$ tels que $(T_v/C)(\R)$ soit compact. Pour la topologie usuelle, $U$ est un ouvert.

Pour $v\in U$ et $w\in W(\C)$ tel que $f(w)=v\in V(\R)\subset V(\C)$, il existe un et un seul homomorphisme
\[
h(w):\bS\longrightarrow T_v
\]
tel que $h(w)\circ r=s_w$: après extension des scalaires à $\C$,
\[
h(w)=s_w\circ z\cdot \overline{s_w}\circ\overline z.
\]
Soit $v\in U$. D'après le théorème de conjugaison des sous-groupes compacts maximaux dans $G/C(\R)^\circ$, il existe $g\in G(\R)$, même $g\in G'(\R)^\circ$, tel que
\[
gT_vg^{-1}\subset K_\infty.
\]
Puisque $T_v$ est un tore maximal, $gT_vg^{-1}$ contient le centre de $K_\infty$. On en déduit l'existence de $w\in W(\C)$ tel que $f(w)=v$ et que $h(w)$ soit conjugué à $h$.

On conclut la démonstration de 5.1 en appliquant au diagramme (5.1.1) et à $U$ la variante suivante du théorème d'irréductibilité de Hilbert, dont je dois la démonstration à J.-P. Serre.

\paragraph{Lemme 5.13.} Soient $E$ une extension finie de $\Q$, et un diagramme commutatif
\[
\begin{array}{ccc}
W & \xrightarrow{\ f\ } & V\\
{\scriptstyle p}\downarrow && \downarrow\\
\Spec(E) & \longrightarrow & \Spec(\Q)
\end{array}
\]
dans lequel
\begin{enumerate}[label=(\roman*)]
\item $V$ est un ouvert de Zariski d'un espace affine sur $\Q$;
\item les fibres géométriques de $p$ sont irréductibles, et $W$ est réduit;
\item $f$ est quasi-fini et dominant.
\end{enumerate}
Alors, pour tout ouvert $U\subset V(\R)$, pour la topologie usuelle, et toute extension $F$ de $E$, il existe
\[
v\in V(\Q)\cap U
\]
tel que $f^{-1}(v)$ soit le spectre d'une extension de $E$ linéairement disjointe de $F$.

Soit $W'=W\otimes_EF$. Les hypothèses de 4.12 sont encore vérifiées pour $W'/F$ et
\[
f':W'\longrightarrow W\longrightarrow V,
\]
et il faut construire $v\in V(\Q)\cap U$ tel que
\[
f'^{-1}(v)=f^{-1}(v)\otimes_EF
\]
soit le spectre d'un corps. Ceci nous ramène au cas où $E=F$, qui résulte du théorème d'irréductibilité de Hilbert, cf. Lang, \emph{Diophantine Geometry}, chap. VIII.

\paragraph{Proposition 5.2.} Pour $x\in\MC(G,h)$ et $K$ compact ouvert dans $G(\Af)$, l'image dans ${}_K\MC(G,h)$ de $G(\Af)\cdot x\subset\MC(G,h)$ est dense.

C'est une conséquence de (0.4). Le lecteur vérifiera.

\paragraph{Proposition 5.3.} Soient $E$ un corps muni d'un plongement complexe $\rho$, $X$ un schéma sur $E$ et $Y$ un sous-schéma réduit fermé de $X_{\C}$. Supposons qu'il existe des extensions finies $E_i$ de $E$ dans $\C$, des schémas $M_{i,\alpha}/E_i$ et des $E_i$-morphismes
\[
v_{i,\alpha}:M_{i,\alpha}\longrightarrow X\otimes_EE_i,
\]
soit
\[
v_i:\coprod_\alpha M_{i,\alpha}\longrightarrow X\otimes_EE_i,
\]
tels que $E=\bigcap_i E_i$, que
\[
v_i\Bigl(\bigl(\coprod_\alpha M_{i,\alpha}\bigr)\otimes_{E_i}\C\Bigr)\subset Y
\]
et que
\[
v_i\Bigl(\bigl(\coprod_\alpha M_{i,\alpha}\bigr)\otimes_{E_i}\C\Bigr)
\]
soit dense dans $Y$. Alors $Y$ est défini sur $E$.

\paragraph{Corollaire 5.4.} Soit $u:(G^1,h^1)\to(G^2,h^2)$ comme en 1.14. Soient $M_E(G^1,h^1)$ et $M_E(G^2,h^2)$ des modèles faiblement canoniques des $\MC(G^i,h^i)$ sur un corps de nombre $E$, avec $E(G^i,h^i)\subset E\subset\C$. Alors
\[
v:\MC(G^1,h^1)\longrightarrow\MC(G^2,h^2)
\]
est défini sur $E$, cf. 3.12.

Soient $K^i\subset G^i(\Af)$ des sous-groupes compacts ouverts tels que $v(K^1)\subset K^2$. Prouvons que
\[
v(K^1,K^2):{}_{K^1}\MC(G^1,h^1)\longrightarrow{}_{K^2}\MC(G^2,h^2)
\]
est défini sur $E$. Soit $u:(H,h')\to(G^1,h^1)$ comme en 3.13 et $F=E\cdot E(H,h')$. Pour $g\in G^1(\Af)$, on dispose de $F$-morphismes
\[
\begin{aligned}
a(g):\quad M_F(H,h')&\longrightarrow M_F(G^1,h^1)
 \xrightarrow{\ g\ } M_F(G^1,h^1)
 \longrightarrow {}_{K^1}M_F(G^1,h^1),\\
b(g):\quad M_F(H,h')&\longrightarrow M_F(G^2,h^2)
 \xrightarrow{\ v(g)\ } M_F(G^2,h^2)
 \longrightarrow {}_{K^2}M_F(G^2,h^2),
\end{aligned}
\]
et, après extension des scalaires à $\C$,
\[
v(K^1,K^2)\circ a(g)=b(g).
\]
Grâce à 5.1 et 5.2, on peut appliquer 5.3 à
\[
X={}_{K^1}M_E(G^1,h^1)\times{}_{K^2}M_E(G^2,h^2),
\]
au graphe $Y$ de $v$, aux $E_i=E\cdot E(H,h')$ pour $(H,h',u)$ variable, et aux
\[
(a(g),b(g)):M_{E_i}(H,h')\longrightarrow X_{E_i}.
\]
On trouve que le graphe de $v$ est défini sur $E$.

Dans le cas particulier $(G^1,h^1)=(G^2,h^2)$, $u=\Id$, 5.4 signifie:

\paragraph{Corollaire 5.5.} À isomorphisme unique près, il existe au plus un modèle faiblement canonique de $\MC(G,h)$ sur $E$, pour $E(G,h)\subset E\subset\C$.

\paragraph{Corollaire 5.6.} Si $G'$ est simplement connexe, la loi de réciprocité d'un modèle faiblement canonique de $\MC(G,h)$ sur $E$, pour $E(G,h)\subset E\subset\C$, est donnée par 3.9.1.

Même démonstration que le cas particulier $v:(G,h)\to(T,vh)$ de 5.4.

\paragraph{Corollaire 5.7.} Soit $v:(G^1,h^1)\hookrightarrow(G^2,h^2)$ comme en 1.15. Si $\MC(G^2,h^2)$ admet un modèle faiblement canonique sur $E$, avec
\[
E(G^2,h^2)\subset E(G^1,h^1)\subset E\subset\C,
\]
alors $\MC(G^1,h^1)$ admet un modèle faiblement canonique sur $E$.

Soit $K^i$ un sous-groupe compact ouvert de $G^i(\Af)$. Utilisant 1.15, on se ramène à montrer seulement que si $v(K^1)\subset K^2$ et que
\[
{}_{K^1}\MC(G^1,h^1)\subset{}_{K^2}\MC(G^2,h^2),
\]
alors le sous-schéma ${}_{K^1}\MC(G^1,h^1)$ de
\[
{}_{K^2}\MC(G^2,h^2)={}_{K^2}M_E(G^2,h^2)\otimes_E\C
\]
est défini sur $E$.

Soit $u:(H,h')\to(G^1,h^1)$ comme en 3.13 et, pour $g\in G^1(\Af)$, soit $a(g)$ l'homomorphisme composé, défini sur $F=E\cdot E(H,h')$:
\[
a(g):M_F(H,h')\longrightarrow M_F(G^1,h^1)
 \xrightarrow{\ v(g)\ } M_F(G^2,h^2)
 \longrightarrow{}_{K^2}M_F(G^2,h^2).
\]
Grâce à 5.1 et 5.2, on conclut en appliquant 5.3 à
\[
X={}_{K^2}M_E(G^2,h^2),\qquad
Y={}_{K^1}\MC(G^1,h^1),
\]
aux $E_i=E\cdot E(H,h')$ pour $(H,h',u)$ variable et aux
\[
a(g):M_{E_i}(H,h')\longrightarrow X_{E_i}.
\]

\paragraph{Exemple 5.8.} Pour $(G,h_0)$ comme en 4.9, on construit par 5.7 et 4.21 un modèle canonique
\[
M(G^\circ,h_0)\subset M(\Gp(V,\psi),h_0)\otimes_{\Q}E(G^\circ,h^\circ).
\]

\paragraph{Variante 5.9.} Soit $(G_1,h_0)$ tel que, avec les notations de 4.22, il existe
\[
u:(G_1,h_0)\hookrightarrow(\Gp_{2g}^{(F)},h_0),
\]
par exemple $(G_1^\circ,h_0)$ avec $(G_1,h_0)$ comme en 4.13. Une application de 5.7 fournit un modèle canonique de $M(G_1,h_0)$.

\paragraph{Proposition 5.10.} Soient $(G,h)$ comme en 3.13, des corps de nombres
\[
E_i:\ E(G,h)\subset E_i\subset\C,
\]
$E$ l'intersection des $E_i$, et $M_{E_i}(G,h)$ un modèle faiblement canonique de $\MC(G,h)$ sur $E_i$. Il existe alors un unique modèle $M_E(G,h)$ de $\MC(G,h)$ sur $E$, tel que pour tout $i$ le modèle $M_E(G,h)\otimes_EE_i$ sur $E_i$ soit isomorphe à $M_{E_i}(G,h)$.

Vu l'assertion d'unicité, on peut supposer les $E_i$ en nombre fini. Soit $F$ l'extension de $E$ dans $\C$ qu'ils engendrent, et soit $\mathfrak E$ l'ensemble des sous-extensions de $F$ contenant l'un des $E_i$. Utilisant 5.5, on obtient pour tout $F'\in\mathfrak E$ un modèle $M_{F'}(G,h)$, et pour toute inclusion $F'\subset F''$, $F',F''\in\mathfrak E$, un isomorphisme de modèles
\[
\alpha(F'',F'):M_{F'}(G,h)\otimes_{F'}F''\longrightarrow M_{F''}(G,h).
\]

Soit $u:(H,h')\to(G,h)$ comme 3.13, avec $E(H,h')$ linéairement disjoint de $F$ sur $E(G,h)$, cf. 5.1. Posons
\[
E_1=E(H,h')\cdot E,\qquad F'_1=E_1\cdot F'=E(H,h')\cdot F'\qquad(F'\in\mathfrak E).
\]
Soit $g\in G(\Af)$ et $K\subset G(\Af)$ un sous-groupe compact ouvert. Puisque $M_{F'}(G,h)$ est faiblement canonique, on dispose d'un $F'_1$-morphisme
\[
a'(g):M_{F'_1}(H,h')\longrightarrow M_{F'_1}(G,h)
 \xrightarrow{\ g\ } M_{F'_1}(G,h)
 \longrightarrow {}_K M_{F'_1}(G,h).
\]
Puisque $F'_1=E_1\otimes_EF'$, on peut « l'interpréter » comme un morphisme de $F'$-schémas
\[
a(g):M_{E_1}(H,h')\otimes_EF'\longrightarrow {}_K M_{F'}(G,h).
\]
Les $a(g)$ vérifient
\[
\alpha(F'',F')\,a(g)_{F''}=a(g).
\]
On conclut par 5.2 et le lemme suivant, appliqué à
\[
X^1=\coprod_{g\in G(\Af)}M_{E_1}(H,h').
\]
Sa démonstration est laissée en exercice au lecteur.

\paragraph{Lemme 5.10.1.} Soient $F/E$ une extension finie de corps et $\mathfrak E$ un ensemble de sous-corps de $F$, d'intersection réduite à $E$, telle que
\[
F\in\mathfrak E,\qquad F'\subset F''\subset F\Longrightarrow F''\in\mathfrak E.
\]
\begin{enumerate}[label=(\roman*)]
\item Le foncteur qui à un schéma $X/E$ associe le système $\mathcal X$ des schémas $X_{F'}/F'$, pour $F'\in\mathfrak E$, et des isomorphismes
\[
X_{F'}\otimes_{F'}F''\xrightarrow{\sim}X_{F''}\qquad(F'\subset F'',\ F',F''\in\mathfrak E)
\]
est pleinement fidèle.
\item Pour qu'un système
\[
\mathcal X=\{X(F')/F'\ (F'\in\mathfrak E);\ \alpha(F'',F'):X(F')\otimes_{F'}F''\xrightarrow{\sim}X(F'')\}
\]
provienne d'un schéma $X/E$, il suffit que les $X(F')/F'$ soient quasi-projectifs et qu'il existe un schéma $X^1/E$ et un morphisme
\[
u:X^1_{\mathfrak E}\longrightarrow\mathcal X
\]
tel que les $u(F'):X^1_{F'}\to X(F')$ soient d'image dense.
\end{enumerate}

\paragraph{Remarque 5.10.2.} Sous les hypothèses de 5.10, pour $u:(H,h')\hookrightarrow(G,h)$ comme en 3.13, le corps de définition de
\[
u:\MC(H,h')\longrightarrow\MC(G,h)
\]
est contenu dans l'extension
\[
\bigcap_i E(H,h')\cdot E_i
\]
de $E(H,h')\cdot E$.

\paragraph{Proposition 5.11.} Soient $h:\bS\to G_\R$ comme en 3.13 et $\delta:\bS\to C^\circ_\R$, où $C$ est le centre de $G$. Si
\[
E(G,h),\ E(C,\delta)\subset E\subset\C
\]
et si $\MC(G,h)$ admet un modèle faiblement canonique sur $E$, alors $\MC(G,h\delta)$ en admet un aussi.

Soit $K$ un sous-groupe compact ouvert de $G(\Af)$ et $L=K\cap C^\circ(\Af)$. Sur $\C$, le morphisme
\[
u:(G\times C^\circ,(h,\delta))\longrightarrow(G,h\delta),\qquad (g,c)\longmapsto gc,
\]
induit des morphismes
\[
u_K:{}_K\MC(G,h)\times {}_L\MC(C^\circ,\delta)
 \longrightarrow {}_K\MC(G,h\delta).
\]
Soit $d$ le morphisme
\[
d:C^\circ\longrightarrow G\times C^\circ,\qquad c\longmapsto(c,c^{-1}).
\]
Le sous-groupe $d(C^\circ(\Af))$ de $G\times C^\circ(\Af)$ agit sur
\[
\MC(G,h)\times\MC(C^\circ,\delta)
\]
via le quotient $\pi_0(C^\circ(\A))$ de $C^\circ(\Af)$, et $u_K$ induit un isomorphisme
\[
\bar u_K:\bigl({}_K\MC(G,h)\times{}_L\MC(C^\circ,\delta)\bigr)/
\bigl(L\backslash\pi_0(C^\circ(\A))\bigr)
\xrightarrow{\sim}{}_K\MC(G,h\delta).
\]
On définit ${}_K M_E(G,h\delta)$ comme le quotient
\[
\bigl({}_K M_E(G,h)\times{}_L M_E(C^\circ,\delta)\bigr)/
\bigl(L\backslash\pi_0(C^\circ(\A))\bigr)
\]
et on vérifie qu'on trouve ainsi un modèle faiblement canonique. Par construction, on dispose de
\[
\tag{5.11.1}
M_E(G,h)\times_E M_E(C^\circ,\delta)\longrightarrow M_E(G,h\delta).
\]

\paragraph{Conjecture 5.12.} Pour tout $(G,h_0)$ comme 3.13, il existe un modèle canonique de $\MC(G,h_0)$.

Supposons, c'est un cas typique, que $G$ soit $\Q$-simple et adjoint. Soit $\widetilde G$ le revêtement universel de $G$ et $\omega$ un poids de $\widetilde G_{\C}$. Si, avec les notations de 3.7, la classe de conjugaison de $h_0r$ est décrite par des entiers $(n_i)_{i\in|\Delta|}$, et que
\[
\omega=\sum m_i\omega_i,
\]
on pose
\[
\langle\omega,h_0r\rangle=\sum n_i m_i.
\]
Jusqu'ici, on n'a pu attaquer 5.12 que lorsque $\widetilde G$ admet une représentation non triviale $V$ telle que toutes les composantes irréductibles $V'$ de $V_{\C}$ vérifient:
\[
(*)\qquad
\text{pour }\omega\text{ parcourant les poids de }V',\ 
\langle\omega,h_0r\rangle\text{ prend au plus deux valeurs.}
\]
Soit $\sigma$ l'involution d'opposition. La condition $(*)$ signifie encore que le poids dominant $\omega$ de $V'$ vérifie
\[
\langle\omega,h_0r\rangle+\langle\sigma\omega,h_0r\rangle=0\text{ ou }1.
\]
Il n'existe pas de telle représentation $V$ si $G$ est exceptionnel, les facteurs non compacts de $G_{\R}$ sont alors de type $E_{6(-14)}$ ou $E_{7(-25)}$, ou de type $D_4$ trialitaire, ni si $G$ est de type $D_n$ et que $G_{\R}$ a des facteurs non compacts des deux types $D_n^{\mathbb H}$ et $D_n^{\mathbb R,2}$.

\section*{\S 6. Modèles étranges}

\paragraph{6.1.} Soit $F$ un corps de nombres totalement réel, et soit
\[
\tag{6.1.1}
{}_F\Sigma:\quad 0\longrightarrow {}_FG'\longrightarrow {}_FG
 \longrightarrow \Gm\longrightarrow0
\]
une forme sur $F$ de la suite exacte
\[
\tag{6.1.2}
{}_F\Sigma_0:\quad 0\longrightarrow\Sp_{2n}\longrightarrow\Gp_{2n}
 \xrightarrow{\ \mu\ }\Gm\longrightarrow0,
\]
où $\Gp$ désigne un groupe de similitudes symplectiques.

On suppose que, pour chaque place réelle $\tau$ de $F$, $\Sigma\otimes_{F,\tau}\R$ est de l'un des types suivants:
\begin{enumerate}[label=(\alph*)]
\item ${}_F\Sigma\otimes_{F,\tau}\R$ est isomorphe à ${}_F\Sigma_0\otimes_{F,\tau}\R$, ou
\item $G'\otimes_{F,\tau}\R$ est compact.
\end{enumerate}
On désigne par $\tau_1,\ldots,\tau_g$ les places réelles de $F$, et on suppose que les places de type (a) sont les places $\tau_1,\ldots,\tau_r$, avec $r>0$.

\paragraph{6.2.} Le groupe ${}_FG$ peut se décrire de la manière suivante:
\begin{enumerate}[label=(\alph*)]
\item On choisit une algèbre de quaternions $B$ sur $F$, indéfinie en les places $\tau_i$ pour $i\le r$, et définie pour $i>r$. On désigne par $x\mapsto\overline{x}$ son involution canonique.
\item On choisit un $B$-module libre $V$, muni d'une forme $F$-bilinéaire non dégénérée $\Phi$, symétrique et telle que
\[
\Phi(bx,y)=\Phi(x,\overline b\,y).
\]
\item On suppose que, pour $i>r$, la forme déduite de $\Phi$ par l'extension des scalaires $\tau_i:F\to\R$ est définie positive.
\item On prend pour ${}_FG$ le groupe des similitudes $B$-linéaires de $\Phi$; en d'autres termes, les $g\in{}_FG(F)$ sont les $g\in\GL_B(V)$ tels qu'il existe $\mu(g)\in F^*$ vérifiant
\[
\Phi(gx,gy)=\mu(g)\Phi(x,y).
\]
\end{enumerate}

\paragraph{6.3.} Désignons par
\[
\Sigma:\quad 0\longrightarrow G'\longrightarrow G\longrightarrow F^*\longrightarrow0
\]
la suite exacte de groupes algébriques sur $\Q$ déduite de ${}_F\Sigma$ par restriction des scalaires à la Weil de $F$ à $\Q$. On a
\[
\tag{6.3.1}
G_{\R}=\prod_{\tau}{}_FG\otimes_{F,\tau}\R
\qquad(\tau\text{ place réelle}).
\]
Un morphisme $h:\bS\to G_\R$ est défini par ses coordonnées
\[
h_i:\bS\longrightarrow {}_FG\otimes_{F,\tau_i}\R.
\]
D'après 1.6, il existe une et une seule classe de conjugaison de morphismes $h$ tels que, pour $i\le r$, $h_i$ soit du type 1.6, et que, pour $i>r$, $h_i$ soit trivial. Soit $h_0$ dans cette classe.

Le corps $E(G,h_0)$ est le sous-corps totalement réel de $\C$ engendré par les
\[
\sum_{i=1}^r\tau_i(f)\qquad(f\in F).
\]

\paragraph{Théorème 6.4.} Avec les notations précédentes, $\MC(G,h_0)$ admet un modèle canonique $M(G,h_0)$.

Soit $Z$ une extension quadratique totalement imaginaire de $F$. Soient
\[
V'=V\otimes_F Z,\qquad L=B\otimes_F Z,
\]
et $G'$ le sous-groupe algébrique de $\GL_L(V')$ engendré par $G$ et $Z^*$:
\[
\tag{6.4.1}
0\longrightarrow F^*\xrightarrow{\ (f,f^{-1})\ }G\times Z^*
 \xrightarrow{\ q\ }G'\longrightarrow0.
\]
Soit $x\mapsto\widetilde x$ l'involution de $L$ produit tensoriel des involutions de $B$ et $Z$. À un facteur dans $F$ près, il existe sur $Z$ une unique forme $F$-bilinéaire alternée $\psi_Z$. Si ${}_F\psi$ est la forme $F$-bilinéaire alternée sur $V'$ produit tensoriel de $\Phi$ et $\psi_Z$, on a
\[
{}_F\psi(\ell x,y)={}_F\psi(x,\widetilde\ell\,y).
\]
Posons
\[
\psi(x,y)=\Tr_{F/\Q}\bigl({}_F\psi(x,y)\bigr).
\]

Choisissons, pour chaque place réelle de $F$, un plongement complexe de $Z$ qui l'induit. On a alors
\[
Z^*(\R)\simeq\prod_{i=1}^g \C^*.
\]
Soit $h_Z$ l'homomorphisme de $\bS$ dans $Z^*_\R$ de coordonnées
\[
h_i:\bS(\R)=\C^*\longrightarrow\C^*\qquad(1\le i\le g)
\]
égales à $z\mapsto1$ pour $i\le r$, à $z\mapsto z^{-1}$ pour $i>r$.

Pour chaque $h:\bS\to G_\R$ comme en 6.3, on pose
\[
h'=q\circ(h\times h_Z):\bS\longrightarrow G'_\R.
\]

\paragraph{Lemme 6.4.2.} Il existe $b\in L$ tel que, pour un choix convenable de $\psi_Z$, la forme
\[
\psi\bigl(x,b\,h'_0(i)y\bigr)\qquad\text{sur }V'\otimes_\Q\R
\]
soit symétrique et définie positive.

Ceci se vérifie après extension des scalaires de $\Q$ à $\R$.

Ce lemme et 5.9 permettent de construire un modèle canonique de $M(G',h'_0)$, car $G'$ est contenu dans le groupe des similitudes symplectiques de ${}_F\psi(x,by)$. Appliquant 5.11 et 5.7, on en déduit un modèle faiblement canonique de $\MC(G,h_0)$ sur l'extension $E(Z^*,h_Z)$ de $E(G,h_0)$. On conclut par 5.10, 5.10.2 et le lemme suivant, qu'on peut déduire de 5.13.

\paragraph{Lemme 6.5.} Soit $F$ un corps de nombre totalement réel, muni d'un ensemble $S$ de plongements réels. Soit $F^\ast$ le corps de nombres correspondant au sous-groupe de $\Gal(\overline{\Q}/\Q)$ qui stabilise $S$. Pour $Z$ une extension quadratique totalement imaginaire de $F$, munie d'un ensemble $T$ de plongements complexes, un au-dessus de chaque élément de $S$, soit de même $Z^\ast\supset F^\ast$ le corps des invariants du stabilisateur de $T$. Alors, pour toute extension $E$ de $F^\ast$, il existe $Z$ et $T$ tels que $Z^\ast$ soit linéairement disjoint de $E$.

\paragraph{Remarque 6.6.} L'application canonique 1.13, 1.14
\[
\tag{6.6.1}
\MC(G,h_0)\times\MC(Z^*,h_Z)\longrightarrow\MC(G',h'_0)
\]
s'interprète de la façon suivante. Chaque conjugué de $h$ définit sur $V_{\C}$ une




bigraduation de Hodge invariante par $F$:
\[
V_{\C}=V^{-1,0}\oplus V^{0,-1}\oplus V^{0,0}
\qquad(\overline{V^{pq}}=V^{qp}).
\]
On prendra garde que la graduation, par le poids, de $V_{\C}$ par les
\[
V^n=\sum_{p+q=n}V^{pq}
\]
n'est pas définie sur $\Q$ si $r\ne g$.

De même, $Z_{\C}$ est bigradué:
\[
Z_{\C}=Z^{-1,0}\oplus Z^{0,-1}\oplus Z^{0,0}.
\]
Le miracle est que le $\Q$-vectoriel $V'=V\otimes_F Z$, muni de la bigraduation de Hodge produit tensoriel
\[
V'_{\C}=V_{\C}^{\prime -1,0}\oplus V_{\C}^{\prime 0,-1}\qquad(!)
\]
est le $H_1$ d'une variété abélienne.

\section*{BIBLIOGRAPHIE}

\begin{description}[leftmargin=3.0em,labelwidth=2.4em,style=nextline]
\item[{[1]}] F. Bruhat et J. Tits -- \emph{Groupes algébriques simples sur un corps local}, Proc. on a conf. on local fields, p. 23--26, Driebergen 1966, Springer-Verlag.
\item[{[2]}] P. Deligne -- \emph{Travaux de Griffiths}, Sém. Bourbaki, 1969/70, exposé n\textsuperscript{o} 376, Lecture Notes 180, Springer-Verlag.
\item[{[3]}] G. Harder -- \emph{Über die Galoiskohomologie halbeinfacher Matrizengruppen II}, Math. Zeitschrift, 92 (1966), p. 396--415.
\item[{[4]}] M. Kneser -- \emph{Schwache Approximation in algebraischen Gruppen}, Coll. sur la théorie des groupes algébriques, p. 41--52, Bruxelles 1962, CBRM.
\item[{[5]}] K. Miyake -- \emph{On models of certain automorphic function fields}.
\item[{[6]}] D. Mumford -- \emph{Geometric invariant theory}, Ergebnisse 34, 1965, Springer-Verlag.
\item[{[7]}] D. Mumford -- \emph{Families of abelian varieties}, in Alg. groups and discontinuous subgroups, p. 347--351, Boulder 1965, AMS.
\item[{[8]}] D. Mumford -- \emph{A Note on Shimura's Paper ``Discontinuous groups and Abelian Varieties''}. Math. Ann. 181, 345--351 (1969).
\item[{[9]}] V. P. Platonov -- \emph{Le problème d'approximation forte et l'hypothèse de Kneser--Tits}, Izvestia Acad. Nauk CCCP, 33 (1969), p. 1211--1219 et 34 (1970), p. 775--777.
\item[{[10]}] G. Shimura -- \emph{On analytic families of polarized abelian varieties and automorphic functions}, Ann. of Math., 78 1 (1963), p. 149--192.
\item[{[11]}] G. Shimura -- \emph{On the field of definition for a field of automorphic functions I, II, III}, Ann. of Math., 80 1 (1964), p. 160--189; 81 1 (1965), p. 124--165 et 83 2 (1966), p. 377--385.
\item[{[12]}] G. Shimura -- \emph{Construction of class fields and zeta functions of algebraic curves}, Ann. of Math. 85 1 (1967), p. 58--159.
\item[{[13]}] G. Shimura -- \emph{Algebraic number fields and symplectic discontinuous groups}, Ann. of Math., 86 3 (1967), p. 503--592.
\item[{[14]}] G. Shimura -- \emph{On canonical models of arithmetic quotients of bounded symmetric domains}, Ann. of Math., 91 1 (1970), p. 144--222.
\item[{[15]}] G. Shimura -- \emph{Automorphic functions and number theory}, Lecture Notes in Math. 54, 1968, Springer-Verlag.
\item[{[16]}] G. Shimura and T. Taniyama -- \emph{Complex multiplication of abelian varieties and its applications to number theory}, Publ. Math. Soc. Jap., 6 (1961).
\end{description}

\end{document}
